\section{Unreduced BV Cotangent, Scalar QME Obstructions, and Cancellations}
\label{sec:app-unreduced-bv-qme}

This appendix computes the scalar part of the mixed brane-defect QME,
the algebraic current presentation available after passing to
distributional principal-part currents, and the precise cancellations
available before scalar reduction or after balanced superization.  Every
statement is formal-local on
\[
  \R^2_{\mathrm{top}}\times\widehat{\C^2}_0,
  \qquad
  \Omega_{\mathrm{loc}}=dz_1\wedge dz_2 .
\]
There is no compact Calabi--Yau target, global BCOV theorem, CoHA,
quintic, OSV/GV, Abel--Jacobi, Igusa, Borcherds, or BKM input in this
appendix.  The calculation separates five formal-local facts: finite
polynomial boundary observables cannot carry the principal-part module;
the current principal-part model carries a strict algebraic \(P_0\)
presentation with zero homotopies inside the added current target; the
ordinary scalar-reduced trace model has a nonzero central CE class; the
unreduced central-character source absorbs the scalar symbol before
scalar reduction; and the standard balanced supertrace model makes the
scalar representative vanish.

\begin{defn}[BV degrees and signs]
\label{def:app-unreduced-bv-degrees}
  The open polynomial BV/Koszul model uses
  \[
    |\phi_1|=|\phi_2|=0,\qquad |\psi=A^*|=-1,\qquad
    Q\psi=[\phi_1,\phi_2],\qquad |Q|=1 .
  \]
  A Hamiltonian generator \(H_f\) has degree \(0\).  A shifted cotangent
  generator \(\eta_\rho=\rho[1]\) is odd of BV degree \(-1\).  The CE
  coordinates satisfy
  \[
    |c_f|=1,\qquad |u_\rho|=0 .
  \]
  This is the cotangent-coordinate grading of the CE/PV block; it is
  not the ordinary graded dual of the BV degree of \(\rho[1]\).  The
  CE/PV cotangent--Schouten bracket has degree \(-1\).  The reduced
  Hamiltonian boundary \(P_0\) bracket has degree \(0\).  The QME
  obstruction differential \(Q+\{I_0,-\}\) has degree \(1\).  The
  Poisson sign convention is the one determined by
  \(\Omega_{\mathrm{loc}}=dz_1\wedge dz_2\), so
  \(\{z_1,z_2\}=1\).  The principal-part coadjoint action is
  \[
    z_1^p z_2^q\cdot\rho_{a,b}
    =
    -(pb-qa+p-q)\rho_{a-p+1,b-q+1}.
  \]
\end{defn}

\begin{prop}[Local geometry of scalar shadows]
\label{prop:app-local-geometry-scalar-shadow}
  All scalar-contact, finite-window, marked-graph, and supertrace
  calculations below use only the formal-local mixed
  holomorphic-topological geometry
  \[
    I\subset\R_{\mathrm{brane}}\subset\R^2_{\mathrm{top}},
    \qquad
    \widehat{\C^2}_0=\operatorname{Spf}\C[[z_1,z_2]],
    \qquad
    \Omega_{\mathrm{loc}}=dz_1\wedge dz_2 .
  \]
  The inverse Poisson bracket is
  \[
    \{f,g\}_{\Omega_{\mathrm{loc}}}
      =
    \partial_{z_1}f\,\partial_{z_2}g
      -
    \partial_{z_2}f\,\partial_{z_1}g .
  \]
  With the convention
  \(\iota_{X_f}\Omega_{\mathrm{loc}}=-df\), one has
  \[
    X_f=(\partial_{z_1}f)\partial_{z_2}
        -(\partial_{z_2}f)\partial_{z_1},
    \qquad
    \operatorname{div}_{\Omega_{\mathrm{loc}}}X_f=0 .
  \]
  Thus no divergence or volume-form anomaly is present in the scalar
  shadow.  The only Hamiltonian scalar cocycle used in this appendix is
  \[
    \omega(f,g)
      =
    [\{f,g\}_{\Omega_{\mathrm{loc}}}]_0,
    \qquad
    \omega(z_1,z_2)=1 .
  \]
  The open index contraction then multiplies this local Hamiltonian scalar
  by the Chan--Paton trace coefficient:
  \[
    \hbar N\,\omega(f,g)\gamma_{\mathbf 1}
    \quad\text{in the ordinary }\lie{gl}_N\text{ branch},
  \]
  \[
    \hbar\,\operatorname{sdim}(\C^{N|N})\,
    \omega(f,g)\gamma_{\mathbf 1}=0
    \quad\text{in the balanced supertrace branch},
  \]
  and, for a marked scalar loop,
  \[
    \hbar\,\operatorname{Str}_{\mathrm{cyc}}(\mathfrak m)\,
    \omega(f,g)\gamma_{\mathbf 1}.
  \]
  A finite window \(M\) only restricts the polynomial Hamiltonian labels,
  graph labels, and counterterm table to a finite local coefficient
  package; it does not change
  \(\Omega_{\mathrm{loc}}\), the Hamiltonian bracket, or the
  divergence-free condition.
\end{prop}

\begin{proof}
  The formula for \(X_f\) is obtained by solving
  \(\iota_{X_f}(dz_1\wedge dz_2)=-df\).  Its divergence is
  \[
    \partial_{z_1}(-\partial_{z_2}f)
      +\partial_{z_2}(\partial_{z_1}f)=0
  \]
  after writing \(X_f=-(\partial_{z_2}f)\partial_{z_1}
  +(\partial_{z_1}f)\partial_{z_2}\).  Therefore the local volume form
  is preserved and cannot contribute an additional scalar contact term.
  Taking the constant Taylor coefficient of the displayed Poisson bracket
  gives the cocycle \(\omega\), with \(\omega(z_1,z_2)=1\).

  The remaining scalar factor is purely the open Chan--Paton index loop.
  The ordinary trace of the identity on \(\C^N\) is \(N\).  The supertrace
  of the identity on \(\C^{N|N}\) is \(N-N=0\).  A marked loop replaces
  the identity by the cyclic product of its marker tensors and gives the
  cyclic supertrace shown above.  Finite-window truncation is a restriction
  of labels and graph tables inside this same formal disk, so it preserves
  the local cocycle and the divergence computation.
\end{proof}

\begin{defn}[Scalar-reduced brane-defect obstruction complex]
\label{def:app-scalar-obstruction-complex}
  Let \(\mathcal O_{\mathrm{loc}}^{\mathrm{bd}}\) denote local
  functionals on the ordinary \(\lie{gl}_N\) brane-defect fields in the
  formal-local model
  \(\R^2_{\mathrm{top}}\times\widehat{\C^2}_0\), with compact support on
  the brane line and with coefficients in the scalar-reduced Hamiltonian
  source \(\bar A=\C[z_1,z_2]/\C\cdot1\).
  The differential controlling a mixed brane-defect QME lift is
  \[
    d_{\mathrm{QME}}=Q+\{I_0,-\}.
  \]
  A scalar contact term is tested by adjoining the degree-one central
  ghost \(\gamma_{\mathbf 1}\) for the omitted constant Hamiltonian
  line.  Its associated graded is the Chevalley--Eilenberg complex of
  \(\bar A\) with trivial coefficients on this central line.
  Writing
  \[
    \mathfrak Q^\bullet_{\mathrm{bd}}
      =(\mathcal O_{\mathrm{loc}}^{\mathrm{bd}},d_{\mathrm{QME}}),
  \]
  the associated-graded scalar-symbol projection is
  \[
    \sigma_{\mathrm{sc}}\colon
    \operatorname{gr}_{\mathrm{sc}}\mathfrak Q^\bullet_{\mathrm{bd}}
    \longrightarrow
    C^\bullet_{\mathrm{Lie}}(\bar A;\C\gamma_{\mathbf 1})[1].
  \]
  A cone or long exact sequence is formed only after one has chosen a
  \(d_{\mathrm{QME}}\)-stable scalar filtration and an actual filtered
  chain projection
  \[
    \widehat\sigma_{\mathrm{sc}}\colon
    \mathfrak Q^\bullet_{\mathrm{bd}}
    \longrightarrow
    C^\bullet_{\mathrm{Lie}}(\bar A;\C\gamma_{\mathbf 1})[1]
  \]
  whose associated graded is the displayed \(\sigma_{\mathrm{sc}}\).
  Without this extra projection, the assertion below is only an
  associated-graded scalar obstruction statement.
  Thus a degree-one QME representative with two Hamiltonian inputs maps to
  a Lie two-cocycle.  Explicitly,
  \[
    \sigma_{\mathrm{sc}}\!\left(
    \kappa(f,g)\int_\R a(t)b(t)\gamma_{\mathbf 1}(t)\,dt\right)
    =
    \kappa(f,g)\gamma_{\mathbf 1}.
  \]
\end{defn}

\begin{lem}[Filtered scalar projection obstruction]\label{lem:filtered-scalar-projection-obstruction}
  Fix a complete decreasing scalar-contact filtration
  \(F^\bullet\mathfrak Q^\bullet_{\mathrm{bd}}\) preserved by
  \(d_{\mathrm{QME}}\), and put the target
  \(C^\bullet_{\mathrm{Lie}}(\bar A;\C\gamma_{\mathbf 1})[1]\) in
  filtration degree \(0\).  Suppose the associated graded carries the
  scalar-symbol projection
  \[
    \sigma_{\mathrm{sc}}\colon
    \operatorname{gr}_{F}\mathfrak Q^\bullet_{\mathrm{bd}}
    \longrightarrow
    C^\bullet_{\mathrm{Lie}}(\bar A;\C\gamma_{\mathbf 1})[1].
  \]
  Choose any filtered linear lift
  \(s\colon\mathfrak Q^\bullet_{\mathrm{bd}}\to
  C^\bullet_{\mathrm{Lie}}(\bar A;\C\gamma_{\mathbf 1})[1]\) of
  \(\sigma_{\mathrm{sc}}\).  Its chain-defect
  \[
    \delta_s=d_{\mathrm{CE}}s-s\,d_{\mathrm{QME}}
  \]
  lowers the scalar-contact filtration.  The first nonzero homogeneous
  piece defines an obstruction class
  \[
    o_\sigma^{(r)}\in
    H^1\!\left(
      \operatorname{Hom}\bigl(\operatorname{gr}_{F}\mathfrak Q^\bullet_{\mathrm{bd}},
      C^\bullet_{\mathrm{Lie}}(\bar A;\C\gamma_{\mathbf 1})[1]\bigr)_{-r}
    \right),
  \]
  with the usual differential
  \(D(\alpha)=d_{\mathrm{CE}}\alpha-(-1)^{|\alpha|}\alpha d_{\operatorname{gr}}\).
  A filtered chain projection
  \(\widehat\sigma_{\mathrm{sc}}\) with associated graded
  \(\sigma_{\mathrm{sc}}\) exists exactly when the obstruction classes
  \(o_\sigma^{(r)}\) vanish successively for all \(r>0\), with compatible
  choices in the completed inverse limit.
\end{lem}

\begin{proof}
  Since the associated graded of \(s\) is already a chain map, the
  commutator \(d_{\mathrm{CE}}s-sd_{\mathrm{QME}}\) strictly lowers the
  filtration.  Its lowest nonzero component is closed for the Hom
  differential because
  \(d_{\mathrm{CE}}^2=d_{\mathrm{QME}}^2=0\).  Replacing \(s\) by
  \(s-\alpha\), where \(\alpha\) lowers filtration by the same amount,
  changes this component by \(D\alpha\).  Thus the first defect can be
  removed exactly when its cohomology class vanishes.  Iterating the same
  argument through the complete filtration gives a compatible sequence of
  corrections if and only if all successive classes vanish.
\end{proof}

\begin{prop}[First scalar-lift obstruction in the ordinary branch]\label{prop:app-first-scalar-lift-obstruction}
  Work in the ordinary scalar-reduced \(\lie{gl}_N\) weighted
  brane-defect local-functional complex
  \(\mathfrak Q^\bullet_{w,\partial,\hbar}(I)\), with \(N>0\).  Let
  \(J_{\mathrm{sc}}\) be the closed ideal of functionals whose
  associated-graded scalar-contact symbol vanishes, and put
  \(F^r=J_{\mathrm{sc}}^r\), completed \(\hbar\)-adically and in the
  finite-window inverse limit.  For any filtered linear lift \(s\) of the
  scalar-symbol projection, the first chain-defect class has Hamiltonian
  two-input evaluation
  \[
    \operatorname{ev}_{f,g}\bigl(o_{\sigma,w}^{(1)}\bigr)
      =
    \hbar N\,\omega(f,g)\gamma_{\mathbf 1}
      \in H^2_{\mathrm{Lie}}(\bar A;\C\gamma_{\mathbf 1})[[\hbar]] .
  \]
  In particular
  \[
    \operatorname{ev}_{z_1,z_2}\bigl(o_{\sigma,w}^{(1)}\bigr)
      =\hbar N\,\gamma_{\mathbf 1}\neq0 .
  \]
  Hence the scalar-symbol projection has no filtered chain lift in the
  ordinary scalar-reduced branch unless one supplies additional data whose
  leading scalar class is \(-\hbar N[\bar c]\), or replaces the source or
  open trace model before scalar reduction.
\end{prop}

\begin{proof}
  The filtration \(F^\bullet\) is the naive scalar-contact filtration:
  its associated graded keeps the leading scalar contact term and
  discards higher scalar-contact corrections.  The first test is the
  bracket of the two linear Hamiltonian boundary insertions.  By
  Theorem~\ref{thm:app-scalar-contact-qme-class},
  \[
    \{I_\partial(a,f),I_\partial(b,g)\}_{\mathrm{open},\hbar}
      -
    I_\partial(ab,\{f,g\}_{\bar A})
      =
    \hbar N\,\omega(f,g)\int_Ia(t)b(t)\,dt
  \]
  after insertion of \(\gamma_{\mathbf 1}\).  Applying the scalar-symbol
  lift gives the displayed representative of
  \(d_{\mathrm{CE}}s-sd_{\mathrm{QME}}\).  For \(f=z_1\) and \(g=z_2\),
  \(\omega(z_1,z_2)=1\).

  If this first defect were a Hom coboundary, evaluation on the
  Hamiltonian two-input subcomplex would give a linear primitive
  \(\bar\eta:\bar A\to\C\gamma_{\mathbf 1}\) with
  \(d_{\mathrm{CE}}\bar\eta=\hbar N\,\omega\gamma_{\mathbf 1}\).  Pulling
  \(\bar\eta\) back to \(A=\C[z_1,z_2]\) and forcing it to vanish on the
  constant line contradicts
  Lemma~\ref{lem:app-unreduced-scalar-primitive}, equivalently
  \(\omega(z_1,z_2)=1\) while \(\{z_1,z_2\}=1\) is zero in \(\bar A\).
  Thus \(o_{\sigma,w}^{(1)}\) is nonzero.

  The obstruction is detected in the finite window containing \(z_1\) and
  \(z_2\).  The admissible weight transports of
  Theorem~\ref{thm:wt-regulator-independence-admissible} identify this
  finite-window class, so changing \(w\) does not affect its vanishing.
  The primitive exists on the unreduced source \(A\), where
  \(\eta(f)=-[f]_0\), and the central-character replacement uses exactly
  this primitive.  In the balanced supertrace branch the displayed
  ordinary-trace evaluation is multiplied by
  \(\operatorname{sdim}\C^{N|N}=0\).  The scalar-contact habitat and
  chain projection constructed in
  Proposition~\ref{prop:app-balanced-scalar-contact-projection} kill this
  obstruction on the balanced scalar-contact branch.  Extension to all
  local functionals remains the separate non-scalar QME obstruction
  problem.
\end{proof}

\begin{thm}[Polynomial unreduced centrality no-go]
\label{thm:app-unreduced-polynomial-centrality-no-go}
  Let \(M_{\mathrm{pol}}\) be any \(Q\)-stable subcomplex of the
  unreduced polynomial open BV observables in which every element has
  finite polynomial degree in \(\phi_1,\phi_2,\psi\).  There is no
  assignment
  \[
    \rho_{a,b}[1]\longmapsto \Theta_{a,b}\in H^\bullet(M_{\mathrm{pol}},Q)
  \]
  and no system of centrality homotopies in \(M_{\mathrm{pol}}\) whose
  induced \(Q\)-cohomology action realizes the principal-part module
  \(\mc P\).  Equivalently, the reduced map
  \(\rho_{a,b}[1]\mapsto\Theta_{a,b}\) cannot be lifted to unreduced
  polynomial boundary representatives with homotopies satisfying
  \[
    \{O_f,\Theta_\rho\}\simeq\Theta_{f\cdot\rho},
    \qquad
    \{\Theta_\rho,\Theta_\sigma\}\simeq0
  \]
  for all polynomial Hamiltonians \(f\).
\end{thm}

\begin{proof}
  The linear Hamiltonians act on polynomial open observables by the
  canonical scalar-translation derivations
  \[
    O_{1,0}\mapsto \partial_{\phi_2},\qquad
    O_{0,1}\mapsto -\partial_{\phi_1}.
  \]
  These derivations commute with \(Q\), since translating
  \(\phi_1\) or \(\phi_2\) by a scalar matrix leaves the moment-map
  equation \(Q\psi=[\phi_1,\phi_2]\) unchanged. On every finite
  polynomial observable, repeated application of either derivation
  eventually vanishes. Hence their induced actions on
  \(H^\bullet(M_{\mathrm{pol}},Q)\) are locally nilpotent.

  On the principal-part module the same two Hamiltonians act by
  \[
    z_1\cdot\rho_{a,b}=-(b+1)\rho_{a,b+1},\qquad
    z_2\cdot\rho_{a,b}=(a+1)\rho_{a+1,b}.
  \]
  Iterating gives
  \[
    z_1^r\cdot\rho_{a,b}
      =(-1)^r(b+1)\cdots(b+r)\rho_{a,b+r}\neq0,
  \]
  and similarly for \(z_2\), for every \(r\geq1\).  Thus no nonzero
  principal part is locally nilpotent for the linear Hamiltonian action.

  If the stated homotopies existed in \(M_{\mathrm{pol}}\), then after
  passing to \(Q\)-cohomology the action of \(O_{1,0}\) and \(O_{0,1}\)
  on the classes \([\Theta_{a,b}]\) would be isomorphic to the displayed
  principal-part action.  Local nilpotence is invariant under module
  isomorphism and under passage to a submodule. This contradicts the
  previous paragraph. Therefore the unreduced lift cannot live in the
  polynomial open BV/Koszul category.
\end{proof}

\begin{lem}[Distributional locality of principal-part currents]
\label{lem:app-distributional-locality-principal-parts}
  Let \(I\subset\R\), \(a\in C^\infty_c(I)\), and
  \(B\in\mathcal D'_c(I)\).  The product
  \[
    (aB)(\varphi)=B(a\varphi),\qquad \varphi\in C^\infty_c(I),
  \]
  is a compactly supported distribution with
  \(\operatorname{supp}(aB)\subset
  \operatorname{supp}(a)\cap\operatorname{supp}(B)\).  Consequently the
  reduced principal-part current brackets
  \[
    \{B_f(a),\Theta_\rho(B)\}
      =\Theta_{f\cdot\rho}(aB),\qquad
    \{\Theta_\rho(B),\Theta_\sigma(C)\}=0
  \]
  are distributionally local on the brane line.  If the supports are
  disjoint, the brackets vanish.
\end{lem}

\begin{proof}
  Multiplication of a compactly supported distribution by a smooth
  function is the standard operation \(B\mapsto aB\) defined by
  precomposition with multiplication by \(a\).  If
  \(\varphi\) is supported away from
  \(\operatorname{supp}(a)\cap\operatorname{supp}(B)\), then
  \(a\varphi\) is supported away from \(\operatorname{supp}(B)\), hence
  \(B(a\varphi)=0\).  This proves the support assertion.  The first
  bracket is the negative transpose of the Hamiltonian action under the
  residue-current pairing, so its interval factor is exactly \(aB\).
  The second bracket is zero because the principal-part current ideal is
  abelian.  Disjoint-support vanishing follows from \(aB=0\).
\end{proof}

\begin{defn}[Central-character principal-part current algebra]
\label{def:app-central-character-principal-part-current-algebra}
  Let \(A=\C[z_1,z_2]\), \(\bar A=A/\C\cdot1\), and let
  \(\chi\colon A\to\C\) be a linear central character on the constant
  Hamiltonian line.  For an interval \(I\subset\R\), write
  \[
    \widetilde{\mathfrak h}_I=\Omega^\bullet_c(I)\widehat\otimes A,
    \qquad
    \mathcal P_I=\mathcal D'_c(I)\widehat\otimes\mc P .
  \]
  The algebraic current target with central character is the completed
  graded-commutative algebra generated by symbols
  \[
    B_f^\chi(a),\qquad \Theta_\rho(B),
    \qquad
    a\in\Omega^0_c(I),\ f\in A,\ B\in\mathcal D'_c(I),\
    \rho\in\mc P,
  \]
  modulo
  \[
    B_{f+c}^\chi(a)=B_f^\chi(a)+\chi(c\cdot1)\int_Ia(t)\,dt,\qquad
    \Theta_{\rho+\sigma}(B)=\Theta_\rho(B)+\Theta_\sigma(B),
  \]
  and with \(P_0\) brackets
  \[
  \begin{aligned}
    \{B_f^\chi(a),B_g^\chi(b)\}
      &=B_{\{f,g\}}^\chi(ab),\\
    \{B_f^\chi(a),\Theta_\rho(B)\}
      &=\Theta_{f\cdot\rho}(aB),\\
    \{\Theta_\rho(B),\Theta_\sigma(C)\}&=0 .
  \end{aligned}
  \]
  Constants act trivially on \(\mc P\).  The product is the ordinary
  graded-commutative factorization product after extension by zero on
  intervals.  This is an algebraic current model; it is not the
  finite polynomial open BV/Koszul complex of
  Theorem~\ref{thm:app-unreduced-polynomial-centrality-no-go}.
\end{defn}

\begin{constr}[Algebraic current presentation map]
\label{constr:app-algebraic-factorization-centre-lift}
  Let
  \[
    \widetilde{\mathfrak g}_I
      =\widetilde{\mathfrak h}_I\ltimes\mathcal P_I[1]
  \]
  with bracket induced by the Hamiltonian bracket on \(A\), the
  coadjoint action on \(\mc P\), and the abelian bracket on
  \(\mathcal P_I[1]\).  In the CE algebra
  \(C^\bullet_{\mathrm{CE},\mathrm{cont}}(\widetilde{\mathfrak g}_I)\),
  write
  \[
    c_{B,\rho}
  \]
  for the degree-one CE coordinate dual to Hamiltonian currents through
  the pairing
  \[
    \langle B\otimes\rho,a(t)dt\otimes f\rangle
      =B(a)\rho(f),
  \]
  and write
  \[
    u_{a(t)dt\otimes f}
  \]
  for the degree-zero CE coordinate dual to the shifted cotangent
  generator under the same pairing with
  \(a(t)dt\otimes f\in\widetilde{\mathfrak h}_I\).
  Define
  \[
    \Phi_I^\chi(c_{B,\rho})=\Theta_\rho(B),\qquad
    \Phi_I^\chi(u_{a(t)dt\otimes f})=B_f^\chi(a),
  \]
  and extend continuously and multiplicatively.
\end{constr}

\begin{thm}[Algebraic current presentation with zero homotopies]
\label{thm:app-algebraic-factorization-centre-zero-homotopies}
  The maps \(\Phi_I^\chi\) of
  Construction~\ref{constr:app-algebraic-factorization-centre-lift}
  define a morphism of \(P_0\) factorization algebras from the
  interval CE source of \(\widetilde{\mathfrak g}_I\) to the
  algebraic current target of
  Definition~\ref{def:app-central-character-principal-part-current-algebra}.
  The factorization centrality homotopies and the \(P_0\)-bracket
  compatibility homotopies in this algebraic model are strict:
  \[
    H_{\mathrm{fact}}=0,\qquad H_{P_0}=0 .
  \]
  Thus the unreduced source with central character constructs the
  \(P_0\) current presentation as far as the distributional algebraic
  current model permits.  This is not constructed as a morphism into
  \(Z^{P_0}_{\mathrm{fact}}(\mathrm{Obs}_{\mathrm{open}})\) for the
  genuine unreduced open BV factorization algebra.  It does not
  contradict the polynomial no-go
  theorem: the image of \(\Theta_\rho(B)\) lives in the added
  principal-part current sector, not in finite polynomial open
  observables.
\end{thm}

\begin{proof}
  The CE differential of
  \(C^\bullet_{\mathrm{CE},\mathrm{cont}}(\widetilde{\mathfrak g}_I)\)
  is dual to exactly three brackets:
  \[
    [\widetilde{\mathfrak h}_I,\widetilde{\mathfrak h}_I],
    \qquad
    \widetilde{\mathfrak h}_I\curvearrowright\mathcal P_I,
    \qquad
    [\mathcal P_I[1],\mathcal P_I[1]]=0 .
  \]
  Under \(\Phi_I^\chi\), these three structure maps become the three
  displayed \(P_0\) brackets of
  Definition~\ref{def:app-central-character-principal-part-current-algebra}.
  Hence the CE differential and the target Schouten differential agree
  on generators, while the degree-zero current \(P_0\) brackets agree
  by the same source brackets.  The biderivation property extends the
  equality to the completed symmetric algebra.

  Factorization compatibility is also strict.  Extension by zero sends
  \(a,B\) to their extensions as test functions and currents; the
  products \(ab\) and \(aB\) are computed after extension.  If the
  supports are disjoint these products vanish, so mixed brackets vanish
  before any homotopy is introduced.  Multiplication in the target is
  graded-commutative, hence multiplication by an image generator is a
  Hochschild central operation with zero product-centrality homotopy.
  The \(P_0\)-centrality defect on generators is the difference between
  the target bracket and the image of the corresponding source bracket;
  it is zero by the previous paragraph.  Therefore
  \(H_{\mathrm{fact}}\) and \(H_{P_0}\) may both be chosen to be zero.

  The same statement holds against arbitrary monomials in the added
  algebraic current target, not only against length-one generators.  Let
  \(Y=Y_1\cdots Y_r\) be a homogeneous source monomial and put
  \(W=\Phi_I^\chi(Y)\), a monomial in the target generators
  \(B_g^\chi(b)\) and \(\Theta_\sigma(C)\).  For a source generator
  \(x\), the product-centrality defect is zero because the target
  product is graded-commutative:
  \[
    \Phi_I^\chi(x)W-(-1)^{|\Phi_I^\chi(x)||W|}W\Phi_I^\chi(x)=0 .
  \]
  Let \(\delta_x\) denote the derivation of the source symmetric
  algebra induced by the corresponding bracket in
  \(\widetilde{\mathfrak g}_I\).  The bracket-compatibility defect
  \[
    D_x(Y)=\{\Phi_I^\chi(x),\Phi_I^\chi(Y)\}
      -\Phi_I^\chi(\delta_xY)
  \]
  is a graded derivation in \(Y\).  It vanishes when \(\Phi_I^\chi(Y)\)
  is \(B_g^\chi(b)\) or \(\Theta_\sigma(C)\) by the three
  displayed brackets.  Hence the Leibniz rule gives \(D_x(Y)=0\) for
  every monomial and, by continuity, for every element of the completed
  algebra.  Thus the centrality homotopies against arbitrary algebraic
  current-target monomials are strict zero homotopies in this model.
  Here ``arbitrary'' means arbitrary monomial in the freely adjoined
  algebraic current target.  It does not mean arbitrary element of the
  genuine unreduced open BV observable factorization algebra, and it
  does not supply the unreduced
  \(Z^{P_0}_{\mathrm{fact}}\)-morphism.
\end{proof}

\begin{lem}[Unreduced scalar primitive]
\label{lem:app-unreduced-scalar-primitive}
  On the local unreduced Hamiltonian algebra \(A=\C[z_1,z_2]\), the
  constant-coefficient cocycle
  \[
    \omega(f,g)=[\{f,g\}_{\Omega_{\mathrm{loc}}}]_0
  \]
  is the CE coboundary of the linear functional
  \[
    \eta(f)=-[f]_0 .
  \]
  With the convention
  \((d_{\mathrm{CE}}\eta)(f,g)=-\eta(\{f,g\})\), one has
  \(d_{\mathrm{CE}}\eta=\omega\).  The primitive does not descend to
  \(\bar A=A/\C\cdot1\).
\end{lem}

\begin{proof}
  The identity is immediate:
  \[
    (d_{\mathrm{CE}}\eta)(f,g)
      =-\eta(\{f,g\}_{\Omega_{\mathrm{loc}}})
      =[\{f,g\}_{\Omega_{\mathrm{loc}}}]_0
      =\omega(f,g).
  \]
  If \(\eta\) descended to \(\bar A\), then it would vanish on
  constants.  But \(\eta(1)=-1\).  Equivalently
  \(\omega(z_1,z_2)=1\) while \(\{z_1,z_2\}=1\) is zero in the
  scalar-reduced Hamiltonian quotient.
\end{proof}

\begin{thm}[Local scalar contact QME obstruction]
\label{thm:app-scalar-contact-qme-class}
  Let
  \[
    \omega(f,g)=[\{f,g\}_{\Omega_{\mathrm{loc}}}]_0
  \]
  be the formal-local Hamiltonian scalar cocycle of
  Proposition~\ref{prop:app-local-geometry-scalar-shadow}.  On monomials,
  \[
    \omega(z_1^kz_2^l,z_1^mz_2^n)=
    \begin{cases}
      kn-lm,&(k+m,l+n)=(1,1),\\
      0,&\text{otherwise.}
    \end{cases}
  \]
  For compactly supported \(a,b\in\Omega^0_c(\R_{\mathrm{brane}})\), the
  unreduced scalar boundary contact in the local mixed HT brane-defect
  model is
  \[
    \{I_\partial(a,f),I_\partial(b,g)\}_{\mathrm{open}}
    =
    I_\partial(ab,\{f,g\}_{\bar A})
    +N\,\omega(f,g)\int_\R a(t)b(t)\,dt .
  \]
  In the quantum Weyl normalization it is equivalently represented by
  \[
    \operatorname{Tr}(A\,XY)-\operatorname{Tr}(A\,YX)
      =
      \hbar N\,\operatorname{Tr}(A)
      \quad\bmod\ \mathcal W_N\widehat\mu(\lie{sl}_N).
  \]
  Equivalently,
  \[
    \operatorname{Tr}(A[X,Y])
      =\hbar N\,\operatorname{Tr}(A)
      \quad\bmod\ \mathcal W_N\widehat\mu(\lie{sl}_N),
  \]
  for the matrix Weyl commutator \([X,Y]=XY-YX\).
  After insertion of the central boundary ghost
  \(\gamma_{\mathbf 1}\), \(|\gamma_{\mathbf 1}|=1\), the degree-one
  obstruction representative is
  \[
    \operatorname{Ob}_{\mathrm{sc}}(\gamma_{\mathbf 1};a,f;b,g)
    =
    \hbar N\,\omega(f,g)
    \int_\R a(t)b(t)\gamma_{\mathbf 1}(t)\,dt .
  \]
  Its associated graded CE class is the formal scalar multiple
  \(\hbar N[\bar c]\in H^2_{\mathrm{Lie}}(\bar A;\C)[[\hbar]]\).  This
  class is not exact in the scalar-reduced Hamiltonian source.  Hence the
  ordinary scalar-reduced \(\lie{gl}_N\) brane-defect model has no
  internal scalar counterterm whose leading symbol cancels this
  obstruction.
\end{thm}

\begin{proof}
  The monomial formula for \(\omega\) follows by applying
  \[
    \{z_1^kz_2^l,z_1^mz_2^n\}
      =(kn-lm)z_1^{k+m-1}z_2^{l+n-1}
  \]
  and then keeping only the constant Taylor coefficient.
  The first displayed formula is
  Corollary~\ref{cor:local-bulk-boundary-coupling}. The second is the
  quantum version of the same central term: the Weyl convention gives
  \([X^i{}_j,Y^k{}_l]=\hbar\delta^i_l\delta^k_j\), and the normal
  ordered moment relation is
  \(Y_\hbar X_\hbar-X_\hbar Y_\hbar+\hbar N I\equiv0\). Hence
  \(XY-YX=\hbar N I\) in the quantum Hamiltonian reduction, and tracing
  against \(A\) gives the stated sign.  The matrix-commutator form is
  the same identity written as
  \(\operatorname{Tr}(A(XY-YX))\).

  The scalar contact itself has degree \(0\). Multiplication by the
  central ghost \(\gamma_{\mathbf 1}\) makes the QME obstruction
  representative degree \(1\), the degree of the obstruction complex.
  The CE associated graded of this representative is the scalar
  extension cocycle \(\hbar N\omega\), hence the formal class
  \(\hbar N[\bar c]\in H^2_{\mathrm{Lie}}(\bar A;\C)[[\hbar]]\).

  If it were exact in the scalar-reduced Hamiltonian source, the leading
  scalar term would give a linear primitive
  \(\bar\eta:\bar A\to\C\) with
  \(\delta\bar\eta=\omega\). Pulling \(\bar\eta\) back to
  \(\C[z_1,z_2]\) and setting the value on \(1\) to zero would contradict
  the computation
  \[
    \omega(z_1,z_2)=1,\qquad \{z_1,z_2\}=1 .
  \]
  Equivalently, Lemma~\ref{lem:app-unreduced-scalar-primitive} gives
  the unreduced scalar primitive used here.  It lives on the unreduced
  polynomial Hamiltonian algebra and does not descend to \(\bar A\).
  Thus the scalar-reduced source has no counterterm whose associated
  graded kills the class.
\end{proof}

\begin{thm}[Central-character source replacement for the scalar QME symbol]
\label{thm:app-central-character-qme-cancellation}
  Let \(A=\C[z_1,z_2]\) and let
  \(\bar f\) denote the image of \(f\) in
  \(\bar A=A/\C\cdot1\) for the local Hamiltonian algebra of
  Proposition~\ref{prop:app-local-geometry-scalar-shadow}.  Define the
  ordinary trace central character
  \[
    \chi_N(f)=N[f]_0
  \]
  and the unreduced boundary generator
  \[
    I^\chi_\partial(a,f)
      =I_\partial(a,\bar f)+\chi_N(f)\int_\R a(t)\,dt .
  \]
  Then, for all \(a,b\in C_c^\infty(\R)\) and all
  \(f,g\in A\),
  \[
    \{I^\chi_\partial(a,f),I^\chi_\partial(b,g)\}_{\mathrm{open}}
      =
    I^\chi_\partial(ab,\{f,g\}).
  \]
  In the quantum Weyl normalization the scalar part is likewise
  absorbed by the same central character:
  \[
    J^\chi_N(1)=N,\qquad
    \operatorname{Tr}(A[X,Y])
      =J^\chi_N([z_1,z_2]_{*})\,\operatorname{Tr}(A)
      =\hbar J^\chi_N(1)\operatorname{Tr}(A).
  \]
  Thus the local scalar QME symbol is absorbed after replacing the
  scalar-reduced source by the unreduced source with central character
  \(\chi_N\).  This is a source replacement before scalar reduction, not a
  null-homotopy inside the original scalar-reduced QME complex.
  Equivalently, if \(\eta(f)=-[f]_0\), then
  \[
    d_{\mathrm{CE}}\eta=\omega,\qquad
    \chi_N(f)=N[f]_0=-N\eta(f),
  \]
  so
  \[
    d_{\mathrm{CE}}\chi_N=-N\omega,\qquad
    \hbar N\omega+d_{\mathrm{CE}}(\hbar\chi_N)=0
  \]
  in the unreduced source.
\end{thm}

\begin{proof}
  Decompose the Poisson bracket into its scalar and reduced parts:
  \[
    \{f,g\}=\overline{\{f,g\}}+\omega(f,g)\cdot1 .
  \]
  By Theorem~\ref{thm:app-scalar-contact-qme-class},
  \[
    \{I_\partial(a,\bar f),I_\partial(b,\bar g)\}_{\mathrm{open}}
      =
    I_\partial(ab,\overline{\{f,g\}})
      +N\omega(f,g)\int_\R a(t)b(t)\,dt .
  \]
  Since the central-character summand is scalar, it has zero open
  bracket with every boundary field.  Therefore the right hand side is
  exactly
  \[
    I_\partial(ab,\overline{\{f,g\}})
    +\chi_N(\{f,g\})\int_\R a(t)b(t)\,dt
    =
    I^\chi_\partial(ab,\{f,g\}).
  \]
  The quantum formula is the same computation in the Weyl algebra:
  \([z_1,z_2]_{*}=\hbar\), \(J^\chi_N(1)=N\), and the matrix moment
  relation gives
  \(\operatorname{Tr}(A[X,Y])=\hbar N\operatorname{Tr}(A)\).  Hence the
  scalar Weyl commutator is the image of the unreduced Moyal scalar
  instead of an obstruction class.

  This is the normal-ordering mechanism.  The relation
  \(YX-XY+\hbar N I\equiv0\) says that an \(XY\)-to-\(YX\) interchange
  contributes the scalar \(\hbar N I\).  In the scalar-reduced source
  this scalar is invisible and therefore appears as
  \(\operatorname{Ob}_{\mathrm{sc}}\).  In the unreduced source with
  \(J_N^\chi(1)=N\), the same scalar is
  \(J_N^\chi([z_1,z_2]_*)=\hbar N\), so it is absorbed before forming
  the QME obstruction class.
\end{proof}

\begin{thm}[Local \(\mathfrak{gl}(N|N)\) supertrace scalar QME cancellation]
\label{thm:app-balanced-supertrace-qme-cancellation}
  Let \(V=\C^{N|N}\), replace the ordinary trace by the supertrace, and
  use the canonical even Weyl relations on \(\operatorname{End}(V)\)
  with the supertrace pairing in the same formal-local model of
  Proposition~\ref{prop:app-local-geometry-scalar-shadow}.  Put
  \[
    \operatorname{sdim}V=\operatorname{Str}(\operatorname{id}_V)=0 .
  \]
  The classical scalar contact term is
  \[
    \{I^{\operatorname{str}}_\partial(a,f),
      I^{\operatorname{str}}_\partial(b,g)\}_{\mathrm{open}}
    =
    I^{\operatorname{str}}_\partial(ab,\{f,g\}_{\bar A})
    +\operatorname{sdim}(V)\,
      \omega(f,g)\int_\R a(t)b(t)\,dt,
  \]
  and hence has zero scalar component for \(V=\C^{N|N}\).  The quantum
  scalar commutator is
  \[
    \operatorname{Str}(A[X,Y])
      =
    \hbar\,\operatorname{sdim}(V)\operatorname{Str}(A)
      =0
      \quad\bmod\ \mathcal W_{\operatorname{End}(V)}
             \widehat\mu(\lie{gl}(V)).
  \]
  Consequently the scalar part of the local mixed brane-defect QME
  vanishes in the balanced \(\lie{gl}(N|N)\) supertrace model without
  adding a central-character counterterm.
\end{thm}

\begin{proof}
  The scalar term in the classical boundary bracket is the contraction
  of the identity endomorphism in the trace pairing.  For an ordinary
  \(N\)-dimensional vector space this contraction is \(N\).  For a
  super vector space it is
  \(\operatorname{Str}(\operatorname{id}_V)=\operatorname{sdim}V\).
  The same local calculation as in
  Theorem~\ref{thm:app-scalar-contact-qme-class} therefore gives the
  displayed classical formula.  For \(V=\C^{N|N}\), this coefficient is
  zero.

  In coordinates, the super-Weyl contraction carries the Koszul sign of
  the supertrace pairing.  Summing the diagonal contraction gives
  \[
    YX-XY+\hbar\,\operatorname{sdim}(V)I\equiv0
      \pmod{\mathcal W_{\operatorname{End}(V)}
             \widehat\mu(\lie{gl}(V))}.
  \]
  Thus
  \[
    \operatorname{Str}(A[X,Y])
      =\hbar\,\operatorname{sdim}(V)\operatorname{Str}(A),
  \]
  which is zero for \(V=\C^{N|N}\).  These are precisely the scalar
  terms entering the degree-one representative
  \(\operatorname{Ob}_{\mathrm{sc}}\); hence that representative is
  zero in the balanced supertrace theory.
\end{proof}

\begin{defn}[Balanced scalar-contact habitat]
\label{def:app-balanced-scalar-contact-habitat}
  Let \(V=\C^{N|N}\).  A balanced scalar-contact habitat over an interval
  \(I\subset\R_{\mathrm{brane}}\) is a complete filtered subcomplex
  \[
    \mathfrak H^\bullet_{\mathrm{bal,sc}}(I)
      \subset
    \mathfrak Q^\bullet_{w,\partial,\hbar,\operatorname{str}}(I)
  \]
  of the local supertrace brane-defect local-functional complex, with
  differential \(d_{\mathrm{QME}}\), satisfying the following conditions.
  First, the filtration \(F^\bullet\) is the closed scalar-contact
  filtration generated by identity-endomorphism contractions in
  \(\operatorname{End}(V)\).  Second, the associated-graded scalar-contact
  extractor factors through the superdimension:
  \[
    \sigma_{\mathrm{bal,sc}}
      =
    \operatorname{sdim}(V)\,\tau_{\mathrm{sc}}
      \colon
    \operatorname{gr}_{F}\mathfrak H^\bullet_{\mathrm{bal,sc}}(I)
      \longrightarrow
    C^\bullet_{\mathrm{Lie}}(\bar A;\C\gamma_{\mathbf 1})[1][[\hbar]]
  \]
  for a continuous \(\C[[\hbar]]\)-linear extractor
  \(\tau_{\mathrm{sc}}\).  Third, the same factorization holds after
  applying \(d_{\mathrm{QME}}\): no scalar-contact symbol is
  present in this habitat except those obtained by contracting a single
  identity endomorphism in the supertrace pairing.
\end{defn}

\begin{prop}[Balanced scalar-contact projection]
\label{prop:app-balanced-scalar-contact-projection}
  On every balanced scalar-contact habitat of
  Definition~\ref{def:app-balanced-scalar-contact-habitat} with
  \(V=\C^{N|N}\), the scalar-symbol map is the zero associated-graded map,
  and it admits the filtered chain projection
  \[
    \widehat\sigma_{\mathrm{bal,sc}}=0\colon
    \mathfrak H^\bullet_{\mathrm{bal,sc}}(I)
      \longrightarrow
    C^\bullet_{\mathrm{Lie}}(\bar A;\C\gamma_{\mathbf 1})[1][[\hbar]] .
  \]
  Its chain-map defect is
  \[
    \delta_{\mathrm{bal}}
      =
    d_{\mathrm{CE}}\widehat\sigma_{\mathrm{bal,sc}}
      -\widehat\sigma_{\mathrm{bal,sc}}d_{\mathrm{QME}}
      =0.
  \]
  Hence all scalar-projection obstruction classes
  \(o^{(r)}_{\sigma,w,\mathrm{bal}}(I)\) vanish on this habitat.  In
  particular the first two-Hamiltonian-input defect is
  \[
    \operatorname{ev}_{f,g}
      \bigl(o^{(1)}_{\sigma,w,\mathrm{bal}}(I)\bigr)
      =
    \hbar\,\operatorname{sdim}(V)\,
    \omega(f,g)\gamma_{\mathbf 1}
      =0 .
  \]
\end{prop}

\begin{proof}
  By definition of the habitat, every scalar-contact symbol factors as
  the ordinary scalar contact multiplied by the contraction of the
  identity endomorphism in the supertrace pairing.  That contraction is
  \[
    \operatorname{Str}(\operatorname{id}_V)
      =
    \operatorname{sdim}(V)=N-N=0 .
  \]
  Therefore
  \(\sigma_{\mathrm{bal,sc}}
    =\operatorname{sdim}(V)\tau_{\mathrm{sc}}=0\) on the associated
  graded.  We choose the zero filtered lift to the stated scalar CE
  target.  Since both composites
  \(d_{\mathrm{CE}}\widehat\sigma_{\mathrm{bal,sc}}\) and
  \(\widehat\sigma_{\mathrm{bal,sc}}d_{\mathrm{QME}}\) are zero, the
  chain-map defect is zero.  Lemma~\ref{lem:filtered-scalar-projection-obstruction}
  then gives \(o^{(r)}_{\sigma,w,\mathrm{bal}}(I)=0\) for every
  \(r>0\).

  Evaluating on two Hamiltonian boundary insertions gives the same local
  contraction as Theorem~\ref{thm:app-balanced-supertrace-qme-cancellation}:
  the ordinary coefficient \(N\) in
  \(\hbar N\,\omega(f,g)\gamma_{\mathbf 1}\) is replaced by
  \(\operatorname{sdim}(V)\).  Thus
  \[
    \hbar\,\operatorname{sdim}(V)\,
    \omega(f,g)\gamma_{\mathbf 1}=0,
  \]
  including the test \((f,g)=(z_1,z_2)\).
\end{proof}

\begin{rmk}[Extension obstruction outside the balanced habitat]
\label{rmk:app-balanced-projection-extension-obstruction}
  The preceding proposition is a positive scalar-contact projection
  theorem, not a full non-scalar graph/QME theorem.  If one enlarges
  \(\mathfrak H^\bullet_{\mathrm{bal,sc}}(I)\) to a complete filtered
  subcomplex \(\mathfrak H'\) carrying scalar symbols not known to factor
  through \(\operatorname{sdim}(V)\), the obstruction is exactly the
  first nonzero homogeneous component of
  \[
    d_{\mathrm{CE}}s-sd_{\mathrm{QME}}
      \in
    \operatorname{Hom}\!\left(
      \operatorname{gr}_{F}\mathfrak H',
      C^\bullet_{\mathrm{Lie}}(\bar A;\C\gamma_{\mathbf 1})[1][[\hbar]]
    \right)
  \]
  for a filtered scalar-symbol lift \(s\).  Its cohomology class is the
  \(o_\sigma^{(r)}\)-class of
  Lemma~\ref{lem:filtered-scalar-projection-obstruction}.  Thus an
  enlargement is admitted exactly when these classes vanish successively,
  with compatible finite-window choices.
\end{rmk}

\begin{defn}[Balanced Costello graph/counterterm subhabitat]
\label{def:app-balanced-costello-graph-subhabitat}
  Fix a Hamiltonian brane-defect specialization of the finite-scale
  Costello graph calculus and a finite-window tower for the formal-local
  mixed HT theory
  \(\R^2_{\mathrm{top}}\times\widehat{\C^2}_0\).  A Costello
  graph/counterterm enlargement
  \[
    \mathfrak H^\bullet_{\mathrm{bal,sc}}(I)
      \subset
    \mathfrak G^\bullet_{\mathrm{bal,Cost}}(I)
      \subset
    \mathfrak Q^\bullet_{w,\partial,\hbar,\operatorname{str}}(I)
  \]
  is a balanced graph subhabitat if it is a complete
  \(d_{\mathrm{QME}}\)-stable filtered subcomplex generated, over
  \(\mathfrak H^\bullet_{\mathrm{bal,sc}}(I)\), by finite-window
  Costello graph contributions \(G_{\Gamma,L,M}\) and local counterterm
  representatives \(C_{\Gamma,L,M}\) such that every scalar-contact
  associated-graded term is an identity-loop supertrace contraction.
  Equivalently, for each homogeneous scalar-contact piece,
  \[
    \sigma_{\mathrm{sc}}\bigl(\operatorname{gr}_{F}G_{\Gamma,L,M}\bigr)
      =
    \operatorname{Str}(\operatorname{id}_V)\,\tau_{\Gamma,L,M},
    \qquad
    \sigma_{\mathrm{sc}}\bigl(\operatorname{gr}_{F}C_{\Gamma,L,M}\bigr)
      =
    \operatorname{Str}(\operatorname{id}_V)\,\tau^C_{\Gamma,L,M},
  \]
  for continuous scalar extractors \(\tau_{\Gamma,L,M}\) and
  \(\tau^C_{\Gamma,L,M}\), and the same factorization holds after
  applying \(d_{\mathrm{QME}}\) to these graph and counterterm
  generators.  The factorization is required to commute with
  finite-window truncation and admissible weight transport.
\end{defn}

\begin{prop}[Balanced graph/counterterm preservation]
\label{prop:app-balanced-costello-graph-preservation}
  Let \(V=\C^{N|N}\).  Every balanced graph subhabitat
  \(\mathfrak G^\bullet_{\mathrm{bal,Cost}}(I)\) of
  Definition~\ref{def:app-balanced-costello-graph-subhabitat} is a
  balanced scalar-contact habitat in the sense of
  Definition~\ref{def:app-balanced-scalar-contact-habitat}.  Hence the
  zero projection
  \[
    \widehat\sigma_{\mathrm{bal,Cost}}=0\colon
    \mathfrak G^\bullet_{\mathrm{bal,Cost}}(I)
      \longrightarrow
    C^\bullet_{\mathrm{Lie}}(\bar A;\C\gamma_{\mathbf 1})[1][[\hbar]]
  \]
  is a filtered chain map, and all scalar-projection obstruction classes
  vanish on this specified graph/counterterm subcomplex.
\end{prop}

\begin{proof}
  By Definition~\ref{def:app-balanced-costello-graph-subhabitat}, each
  new graph or counterterm scalar symbol is obtained by closing an
  identity endomorphism loop in the supertrace pairing.  Its coefficient
  is therefore
  \[
    \operatorname{Str}(\operatorname{id}_V)
      =
    \operatorname{sdim}(V)=N-N=0.
  \]
  Thus the associated-graded scalar extractor is zero on the graph and
  counterterm generators, and it is already zero on
  \(\mathfrak H^\bullet_{\mathrm{bal,sc}}(I)\) by
  Proposition~\ref{prop:app-balanced-scalar-contact-projection}.  The
  same factorization after applying \(d_{\mathrm{QME}}\) says that the
  differential does not create a new scalar-contact symbol outside the
  identity-loop class.  Therefore
  \[
    d_{\mathrm{CE}}\widehat\sigma_{\mathrm{bal,Cost}}
      -
    \widehat\sigma_{\mathrm{bal,Cost}}d_{\mathrm{QME}}=0.
  \]
  The finite-window and weight-transport compatibility in the definition
  gives the same zero class in the completed inverse limit.  Hence the
  scalar-contact projection obstruction tower of
  Lemma~\ref{lem:filtered-scalar-projection-obstruction} vanishes on
  \(\mathfrak G^\bullet_{\mathrm{bal,Cost}}(I)\).
\end{proof}

\begin{thm}[Graph-extension scalar obstruction]
\label{thm:app-costello-graph-extension-obstruction}
  Let
  \[
    \mathfrak H^\bullet_{\mathrm{bal,sc}}(I)
      \subset
    \mathfrak G^\bullet_{\mathrm{Cost}}(I)
      \subset
    \mathfrak Q^\bullet_{w,\partial,\hbar,\operatorname{str}}(I)
  \]
  be any complete \(d_{\mathrm{QME}}\)-stable filtered
  Costello graph/counterterm enlargement.  Put
  \[
    \overline{\mathfrak G}^\bullet_{\mathrm{Cost}}(I)
      =
    \mathfrak G^\bullet_{\mathrm{Cost}}(I)/
    \mathfrak H^\bullet_{\mathrm{bal,sc}}(I),
    \qquad
    S^\bullet=
    C^\bullet_{\mathrm{Lie}}(\bar A;\C\gamma_{\mathbf 1})[1][[\hbar]].
  \]
  Let
  \[
    \bar\sigma_{\mathrm{Cost}}\colon
    \operatorname{gr}_{F}
      \overline{\mathfrak G}^\bullet_{\mathrm{Cost}}(I)
      \longrightarrow S^\bullet
  \]
  be the associated-graded scalar-symbol extractor of the added
  graph/counterterm terms.  The zeroth obstruction to extending the
  balanced scalar-contact projection to
  \(\mathfrak G^\bullet_{\mathrm{Cost}}(I)\) is the actual
  associated-graded chain-map defect
  \[
    \delta^{(0)}_{\mathrm{Cost},\sigma}(I)
      =
      d_{\mathrm{CE}}\bar\sigma_{\mathrm{Cost}}
        -
      \bar\sigma_{\mathrm{Cost}}\,d_{\operatorname{gr}}
    \in
    \operatorname{Hom}^{1}_{0}\bigl(
        \operatorname{gr}_{F}
          \overline{\mathfrak G}^\bullet_{\mathrm{Cost}}(I),
        S^\bullet
      \bigr) .
  \]
  If this defect vanishes, choose a filtered lift \(s\) of
  \(\bar\sigma_{\mathrm{Cost}}\), extended by zero on
  \(\mathfrak H^\bullet_{\mathrm{bal,sc}}(I)\).  If the first nonzero
  homogeneous component of
  \[
    \delta_s=d_{\mathrm{CE}}s-sd_{\mathrm{QME}}
  \]
  lowers the scalar-contact filtration by \(r>0\), then it defines
  \[
    \mathfrak o^{(r)}_{\mathrm{Cost},\sigma}(I)
      =
    \left[\delta_s^{(-r)}\right]
    \in
    H^1\!\left(
      \operatorname{Hom}_{-r}\bigl(
        \operatorname{gr}_{F}
          \overline{\mathfrak G}^\bullet_{\mathrm{Cost}}(I),
        S^\bullet
      \bigr)
    \right).
  \]
  These classes are independent of the chosen lift up to the usual Hom
  coboundary.  A filtered chain projection extending the balanced
  projection exists exactly when
  \[
    \delta^{(0)}_{\mathrm{Cost},\sigma}(I)=0,\qquad
    \mathfrak o^{(r)}_{\mathrm{Cost},\sigma}(I)=0
    \quad(r>0)
  \]
  successively, with compatible finite-window representatives.

  Explicitly, if \(X_{\Gamma,M}\) is an added graph contribution or
  counterterm representative and, modulo \(F^{p-r+1}\),
  \[
    d_{\mathrm{QME}}X_{\Gamma,M}
      =
    \sum_{\Gamma'}\varepsilon_{\Gamma,\Gamma'}X_{\Gamma',M},
    \qquad
    X_{\Gamma,M}\in F^p,
  \]
  with the signs \(\varepsilon_{\Gamma,\Gamma'}\) determined by the
  graph orientation, BV bracket, and counterterm convention, then the
  obstruction cocycle is computed by
  \[
    \mathfrak o^{(r)}_{\mathrm{Cost},\sigma}(X_{\Gamma,M})
      =
    {\left[
      d_{\mathrm{CE}}s(X_{\Gamma,M})
        -
      \sum_{\Gamma'}\varepsilon_{\Gamma,\Gamma'}s(X_{\Gamma',M})
    \right]}_{F^{p-r}/F^{p-r+1}} .
  \]
  Thus a graph or counterterm term outside the identity-loop
  superdimension class is not admitted by the balanced theorem unless
  this displayed relative Hom class is killed by an explicit compatible
  counterterm correction.
\end{thm}

\begin{proof}
  The quotient by
  \(\mathfrak H^\bullet_{\mathrm{bal,sc}}(I)\) makes the statement
  relative to the already proved zero projection on the balanced
  scalar-contact habitat.  If
  \(d_{\mathrm{CE}}\bar\sigma_{\mathrm{Cost}}
    -\bar\sigma_{\mathrm{Cost}}d_{\operatorname{gr}}\)
  is nonzero, then the associated-graded scalar extractor is not a chain
  map; no filtered chain lift can have it as associated graded.  This is
  the zeroth defect
  \(\delta^{(0)}_{\mathrm{Cost},\sigma}(I)\).

  Once the associated-graded defect is killed, the proof is the relative
  version of Lemma~\ref{lem:filtered-scalar-projection-obstruction}.
  The commutator \(d_{\mathrm{CE}}s-sd_{\mathrm{QME}}\) strictly lowers
  filtration.  Its first nonzero homogeneous piece is closed for the Hom
  differential because both differentials square to zero and because the
  lower pieces have already been removed.  Replacing \(s\) by
  \(s-\alpha\), with \(\alpha\) lowering filtration by the same amount
  and vanishing on
  \(\mathfrak H^\bullet_{\mathrm{bal,sc}}(I)\), changes this component
  by the Hom coboundary \(D\alpha\).  Therefore the class is independent
  of choices and vanishes exactly when the next correction to \(s\)
  exists.  Completeness of the filtration and compatibility of the
  finite-window tower give the stated successive criterion.

  The displayed graph formula is this same chain-defect evaluated on a
  homogeneous graph or counterterm generator, with the actual
  \(d_{\mathrm{QME}}\)-boundary written as the signed sum of its graph
  and counterterm boundary terms.  Hence it is computable from the
  supplied Costello graph differential, orientation signs, and scalar
  symbols.  If all added terms satisfy
  Definition~\ref{def:app-balanced-costello-graph-subhabitat}, then
  \(s=0\) is the required lift for \(V=\C^{N|N}\), and all the displayed
  classes vanish, recovering
  Proposition~\ref{prop:app-balanced-costello-graph-preservation}.
\end{proof}

\begin{prop}[First non-identity-loop scalar symbol]
\label{prop:app-first-nonidentity-loop-scalar-symbol}
  Let \(V=\C^{N|N}\).  Fix a finite local Hamiltonian window \(M\)
  containing the polynomial Hamiltonian inputs under discussion, and let
  \(T_\bullet=(T_1,\ldots,T_m)\) be an ordered word of even
  \(\operatorname{End}(V)\)-labels carried by one open index loop.  Write
  \(M(T_\bullet)=T_1\cdots T_m\), with the empty product equal to
  \(\operatorname{id}_V\).  Let
  \(\Gamma^{T_\bullet}_{1,M}\) be the connected local
  one-cross-contraction boundary graph with two Hamiltonian external legs
  \((a,f)\) and \((b,g)\), one Weyl edge, and this marked scalar loop.
  Denote the degree-one scalar-contact representative obtained after
  inserting \(\gamma_{\mathbf 1}\) by
  \(X_{\Gamma^{T_\bullet}_{1,M}}(a,f;b,g)\).

  Then
  \(X_{\Gamma^{T_\bullet}_{1,M}}(a,f;b,g)\in F^1\), its leading term is
  not in \(F^2\) unless the displayed coefficient vanishes, and
  \[
    X_{\Gamma^{T_\bullet}_{1,M}}(a,f;b,g)
      =
    \hbar\,\operatorname{Str}\!\bigl(M(T_\bullet)\bigr)\,
    \omega(f,g)\int_Ia(t)b(t)\gamma_{\mathbf 1}(t)\,dt
      \quad\bmod F^2 .
  \]
  Equivalently,
  \[
    \sigma_{\mathrm{sc}}\!\left(
      \operatorname{gr}_{F}
        X_{\Gamma^{T_\bullet}_{1,M}}(a,f;b,g)\right)
      =
    \hbar\,\operatorname{Str}\!\bigl(M(T_\bullet)\bigr)\,
    \omega(f,g)\gamma_{\mathbf 1}.
  \]
  For monomials,
  \[
    \omega(z_1^kz_2^l,z_1^mz_2^n)=
    \begin{cases}
      kn-lm,&(k+m,l+n)=(1,1),\\
      0,&\text{otherwise,}
    \end{cases}
  \]
  so the first non-identity test is obtained by a single marker \(T\):
  \[
    \sigma_{\mathrm{sc}}\!\left(
      \operatorname{gr}_{F}X_{\Gamma^{T}_{1,M}}(1,z_1;1,z_2)\right)
      =
    \hbar\,\operatorname{Str}(T)\gamma_{\mathbf 1}.
  \]
  In particular, if \(P_+\) is projection onto one even basis line, then
  \[
    \sigma_{\mathrm{sc}}\!\left(
      \operatorname{gr}_{F}X_{\Gamma^{P_+}_{1,M}}(1,z_1;1,z_2)\right)
      =
    \hbar\,\gamma_{\mathbf 1}\neq0.
  \]
  Thus this is the first scalar symbol not killed by the identity-loop
  computation.  Its associated CE class is
  \[
    \hbar\,\operatorname{Str}\!\bigl(M(T_\bullet)\bigr)
    [\bar c]\,\gamma_{\mathbf 1}.
  \]
  It is nonzero exactly when
  \(\operatorname{Str}(M(T_\bullet))\neq0\).  If every scalar-contact
  loop monodromy word in a complete \(d_{\mathrm{QME}}\)-stable
  graph/counterterm subcomplex, and every such word appearing after
  \(d_{\mathrm{QME}}\), has supertrace zero, then the zero scalar
  projection is again a filtered chain map on that subcomplex.
\end{prop}

\begin{proof}
  The graph has one local scalar contact on the brane line, so it lies in
  \(F^1\).  There is no second scalar contact in the one-edge graph, hence
  its associated graded is read in \(F^1/F^2\).  The boundary Wick
  normalization fixed in
  Theorem~\ref{thm:open-line-midpoint-graph-weights} gives the one-edge
  antisymmetrized coefficient \(\hbar P(f,g)\).  Taking the scalar
  Taylor coefficient gives
  \[
    [P(f,g)]_0=\omega(f,g),
  \]
  with positive sign for the ordered pair \((z_1,z_2)\), since
  \(\{z_1,z_2\}=1\).  Reversing the two external legs changes the sign
  through \(\omega(g,f)=-\omega(f,g)\).

  It remains only to compute the open index loop.  In a homogeneous basis
  \(e_\alpha\) of \(V\), the closed marked loop contributes
  \[
    \sum_\alpha (-1)^{|e_\alpha|}
      \bigl(M(T_\bullet)\bigr)^\alpha{}_\alpha
      =
    \operatorname{Str}\!\bigl(M(T_\bullet)\bigr).
  \]
  For the empty word this is
  \(\operatorname{Str}(\operatorname{id}_V)=\operatorname{sdim}V=0\),
  which is exactly the identity-loop cancellation of
  Proposition~\ref{prop:app-balanced-costello-graph-preservation}.  For a
  non-identity marker \(T\) the same contraction gives
  \(\operatorname{Str}(T)\).  Projection onto one even basis line has
  supertrace \(1\), giving the displayed nonzero symbol.

  The scalar contact has BV degree \(0\); multiplication by
  \(\gamma_{\mathbf 1}\), of degree \(1\), puts the representative in the
  degree-one QME obstruction complex.  The cocycle \(\omega\) is the same
  scalar Hamiltonian cocycle as in
  Theorem~\ref{thm:app-scalar-contact-qme-class}.  That theorem proves
  \([\bar c]\neq0\) on the scalar-reduced Hamiltonian source, so multiplying
  by the scalar
  \(\hbar\,\operatorname{Str}(M(T_\bullet))\) gives a nonzero class
  exactly when this supertrace is nonzero.

  Finally suppose the stated marked-loop subcomplex has only
  supertraceless monodromy words, before and after applying
  \(d_{\mathrm{QME}}\).  The formula above then gives zero associated
  scalar symbol on every graph and counterterm generator and also on
  every associated-graded boundary term.  Hence
  \(d_{\mathrm{CE}}0-0\,d_{\mathrm{QME}}=0\), so the zero scalar projection
  is a filtered chain map on that larger, explicitly delimited subcomplex.
\end{proof}

\begin{defn}[Supertraceless-monodromy prehabitat]
\label{def:app-supertraceless-monodromy-prehabitat}
  Let \(V=\C^{N|N}\).  For an ordered word
  \(T_\bullet=(T_1,\ldots,T_m)\) of even \(\operatorname{End}(V)\)-labels
  on a scalar open index loop, set
  \[
    \operatorname{st}(T_\bullet)
      =
    \operatorname{Str}(T_1\cdots T_m),
    \qquad
    \operatorname{st}(\varnothing)
      =
    \operatorname{Str}(\operatorname{id}_V)=0 .
  \]
  The word is supertraceless if \(\operatorname{st}(T_\bullet)=0\).
  This condition is imposed on the ordered product \(T_1\cdots T_m\),
  not on the individual letters.

  Fix a complete \(d_{\mathrm{QME}}\)-stable finite-window
  graph/counterterm enlargement of the local balanced supertrace
  brane-defect complex.  The supertraceless-monodromy
  prehabitat
  \[
    \mathfrak G^{\bullet,\mathrm{pre}}_{\operatorname{st}=0}(I)
  \]
  is the closed filtered span, over
  \(\mathfrak H^\bullet_{\mathrm{bal,sc}}(I)\), of homogeneous
  graph/counterterm representatives whose every scalar-contact
  associated-graded loop word is supertraceless.  Let
  \[
    \Pi_{\operatorname{st}\neq0}
  \]
  denote the quotient of the associated-graded scalar-contact monodromy
  module by the closed span of supertraceless loop words.  The
  \(d_{\mathrm{QME}}\)-preservation obstruction is the family of
  homogeneous quotient classes
  \[
    \mathcal B_{\operatorname{st}=0}(X;q)
      =
    \Pi_{\operatorname{st}\neq0}
    \left([d_{\mathrm{QME}}X]_{F^q/F^{q+1}}\right),
    \qquad
    X\in\mathfrak G^{\bullet,\mathrm{pre}}_{\operatorname{st}=0}(I),
  \]
  over all filtration degrees \(q\).  We write
  \(\mathcal B_{\operatorname{st}=0}=0\) when every class in this family
  vanishes.
  If \(\mathcal B_{\operatorname{st}=0}=0\), the prehabitat is a complete
  \(d_{\mathrm{QME}}\)-stable filtered subcomplex; in that case it is
  called the supertraceless-monodromy graph/counterterm subhabitat and
  is denoted
  \(\mathfrak G^\bullet_{\operatorname{st}=0}(I)\).
\end{defn}

\begin{thm}[Supertraceless-monodromy scalar projection]
\label{thm:app-supertraceless-monodromy-projection}
  Assume the notation of
  Definition~\ref{def:app-supertraceless-monodromy-prehabitat}.  The
  prehabitat is preserved by \(d_{\mathrm{QME}}\) if and only if
  \[
    \mathcal B_{\operatorname{st}=0}=0 .
  \]
  Under this condition the zero scalar projection
  \[
    \widehat\sigma_{\operatorname{st}=0}=0\colon
    \mathfrak G^\bullet_{\operatorname{st}=0}(I)
      \longrightarrow
    C^\bullet_{\mathrm{Lie}}(\bar A;\C\gamma_{\mathbf 1})[1][[\hbar]]
  \]
  is a filtered chain map.  Hence all scalar-projection obstruction
  classes of Lemma~\ref{lem:filtered-scalar-projection-obstruction}
  vanish on the supertraceless-monodromy subhabitat.

  More explicitly, let \(X_{\Gamma,M}\in F^p\) be a homogeneous
  generator of the prehabitat and suppose that, modulo \(F^{p-r+1}\),
  \[
    d_{\mathrm{QME}}X_{\Gamma,M}
      =
    \sum_{\Gamma',U_\bullet}
      \varepsilon_{\Gamma,\Gamma',U_\bullet}
      X_{\Gamma'^{\,U_\bullet},M},
  \]
  where \(U_\bullet\) runs over scalar-loop monodromy words appearing in
  the boundary terms.  The first possible failure of preservation is the
  quotient class
  \[
    \mathcal B^{(r)}_{\operatorname{st}=0}(X_{\Gamma,M})
      =
    \left[
      \sum_{\Gamma',U_\bullet}
      \varepsilon_{\Gamma,\Gamma',U_\bullet}
      X_{\Gamma'^{\,U_\bullet},M}
    \right]_{\operatorname{st}\neq0,\;F^{p-r}/F^{p-r+1}} .
  \]
  Its scalar CE image on a two-Hamiltonian-input scalar loop is
  \[
    \hbar
    \sum_{\Gamma',U_\bullet}
      \varepsilon_{\Gamma,\Gamma',U_\bullet}
      \operatorname{Str}\!\bigl(M(U_\bullet)\bigr)\,
      \omega(f_{\Gamma'},g_{\Gamma'})\gamma_{\mathbf 1}.
  \]
  In cohomology this becomes
  \[
    \hbar
    \sum_{\Gamma',U_\bullet}
      \varepsilon_{\Gamma,\Gamma',U_\bullet}
      \operatorname{Str}\!\bigl(M(U_\bullet)\bigr)\,
      [\bar c]\,\gamma_{\mathbf 1}.
  \]
  Thus a non-supertraceless boundary word is the exact obstruction to
  \(d_{\mathrm{QME}}\)-preservation, and its scalar shadow is precisely
  the marked-loop class computed in
  Proposition~\ref{prop:app-first-nonidentity-loop-scalar-symbol}.

  For the unmarked Hamiltonian boundary generators
  \(I^{\operatorname{str}}_\partial(a,f)\) used in
  Theorem~\ref{thm:app-balanced-supertrace-qme-cancellation}, and for the
  identity-loop graph generators of
  Definition~\ref{def:app-balanced-costello-graph-subhabitat}, the
  obstruction \(\mathcal B_{\operatorname{st}=0}\) vanishes.  These
  generators carry only the empty monodromy word, and the displayed
  differential \(Q+\{I_0,-\}\) contains no fixed noncentral
  \(\operatorname{End}(V)\)-marker.  Therefore applying
  \(d_{\mathrm{QME}}\) to this displayed Hamiltonian generator set cannot
  create a non-supertraceless marked word.  This proves preservation for
  the Hamiltonian generators treated in this appendix; it does not assert
  preservation for a larger Costello specialization with additional
  marked idempotents, holonomies, or counterterm tensors.
\end{thm}

\begin{proof}
  By definition,
  \(\Pi_{\operatorname{st}\neq0}\) kills exactly the associated-graded
  scalar-contact monodromy terms whose loop words are supertraceless.
  Hence
  \(\mathcal B_{\operatorname{st}=0}=0\) says precisely that the
  \(d_{\mathrm{QME}}\)-boundary of every generator of the prehabitat has
  only supertraceless loop words in every associated-graded filtration
  degree.  Since the filtration is complete, this homogeneous condition
  is equivalent to preservation of the closed filtered span.  This proves
  the preservation criterion.

  On the preserved subhabitat every scalar-contact loop word
  \(T_\bullet\) satisfies
  \(\operatorname{Str}(M(T_\bullet))=0\).  Proposition~\ref{prop:app-first-nonidentity-loop-scalar-symbol}
  computes the associated-graded scalar symbol of such a loop as
  \[
    \hbar\,\operatorname{Str}(M(T_\bullet))\,
    \omega(f,g)\gamma_{\mathbf 1},
  \]
  and this is zero.  The same is true after applying
  \(d_{\mathrm{QME}}\), by the preservation criterion.  Therefore both
  composites
  \(d_{\mathrm{CE}}\widehat\sigma_{\operatorname{st}=0}\) and
  \(\widehat\sigma_{\operatorname{st}=0}d_{\mathrm{QME}}\) are zero, so
  the zero projection is a filtered chain map.  Lemma~\ref{lem:filtered-scalar-projection-obstruction}
  then kills the whole scalar-projection obstruction tower on this
  subhabitat.

  The formula for
  \(\mathcal B^{(r)}_{\operatorname{st}=0}(X_{\Gamma,M})\) is just the
  preceding quotient criterion evaluated on a homogeneous generator and
  then reduced modulo \(F^{p-r+1}\).  Applying the one-loop scalar
  computation of
  Proposition~\ref{prop:app-first-nonidentity-loop-scalar-symbol} to each
  boundary term gives the displayed scalar CE representative.  Passing to
  cohomology replaces \(\omega\) by the nonzero class
  \([\bar c]\) from
  Theorem~\ref{thm:app-scalar-contact-qme-class}.

  For the unmarked Hamiltonian boundary generators, the open index loop
  contains no inserted label; its word is \(\varnothing\), so its
  coefficient is
  \(\operatorname{Str}(\operatorname{id}_V)=0\).  The differential
  \(Q+\{I_0,-\}\) appearing in this appendix changes the BV fields and
  takes brackets with the unmarked classical interaction.  It does not
  introduce an external even endomorphism marker on an open scalar loop.
  Thus the image of the unmarked Hamiltonian generator set again has only
  empty loop words.  The identity-loop graph subhabitat has the same
  property by
  Definition~\ref{def:app-balanced-costello-graph-subhabitat}, which
  requires the same factorization after applying \(d_{\mathrm{QME}}\).
  This proves the stated preservation result on the displayed generator
  set and leaves arbitrary marked graph/counterterm tensors governed by
  \(\mathcal B_{\operatorname{st}=0}\).
\end{proof}

\begin{prop}[Marked \(\mathcal B_{\operatorname{st}=0}\) and gauge-equivariant markers]
\label{prop:app-marked-bst0-equivariant-differential}
  Let \(V=\C^{N|N}\), and impose full \(\lie{gl}(V)\)-gauge
  equivariance on open Chan-Paton labels.  A theorem-grade computation
  of \(\mathcal B_{\operatorname{st}=0}\) for a marked
  Costello graph/counterterm enlargement requires, for every homogeneous
  generator \(X_{\Gamma,M}\in F^p\), an actual marked boundary expansion
  \[
    d_{\mathrm{QME}}X_{\Gamma,M}
      =
    \sum_{\Gamma',\lambda}
      \varepsilon_{\Gamma,\Gamma',\lambda}\,
      X_{\Gamma'^{\,\mathfrak m_{\Gamma,\Gamma',\lambda}},M}
      \quad \bmod F^{p-r+1},
    \label{eq:app-marked-differential-datum}
  \]
  where each
  \[
    \mathfrak m_{\Gamma,\Gamma',\lambda}
      \in
    \bigl(\operatorname{End}(V)^{\otimes m_{\lambda}}\bigr)^{\lie{gl}(V)}
  \]
  is an even invariant marker tensor decorating the scalar open index
  loop.  For
  \(\mathfrak m=\sum_\alpha
  T_{\alpha,1}\otimes\cdots\otimes T_{\alpha,m}\), put
  \[
    \operatorname{Str}_{\mathrm{cyc}}(\mathfrak m)
      =
    \sum_\alpha
      \operatorname{Str}(T_{\alpha,1}\cdots T_{\alpha,m}) .
  \]
  Then the first preservation obstruction is
  \[
    \mathcal B^{(r)}_{\operatorname{st}=0}(X_{\Gamma,M})
      =
    \left[
      \sum_{\Gamma',\lambda}
      \varepsilon_{\Gamma,\Gamma',\lambda}\,
      X_{\Gamma'^{\,\mathfrak m_{\Gamma,\Gamma',\lambda}},M}
    \right]_{\operatorname{st}\neq0,\;F^{p-r}/F^{p-r+1}},
  \]
  and its two-Hamiltonian-input scalar shadow is
  \[
    \hbar
    \sum_{\Gamma',\lambda}
      \varepsilon_{\Gamma,\Gamma',\lambda}\,
      \operatorname{Str}_{\mathrm{cyc}}
        (\mathfrak m_{\Gamma,\Gamma',\lambda})\,
      \omega(f_{\Gamma'},g_{\Gamma'})\gamma_{\mathbf 1}.
  \]
  In cohomology this becomes
  \[
    \hbar
    \sum_{\Gamma',\lambda}
      \varepsilon_{\Gamma,\Gamma',\lambda}\,
      \operatorname{Str}_{\mathrm{cyc}}
        (\mathfrak m_{\Gamma,\Gamma',\lambda})\,
      [\bar c]\,\gamma_{\mathbf 1}.
  \]

  On the displayed unmarked Hamiltonian generator set the actual
  differential is \(Q+\{I_0,-\}\), with unmarked classical interaction;
  the balanced graph/counterterm subhabitat only allows identity-loop
  scalar contractions before and after \(d_{\mathrm{QME}}\).  Hence
  \[
    \mathcal B_{\operatorname{st}=0}=0
  \]
  on that displayed surface.  For any larger marked Costello
  specialization not equipped with an oriented finite-window marked
  boundary complex as in
  Definition~\ref{def:app-oriented-full-equivariant-marked-boundary-complex},
  the missing datum blocking a theorem-grade scalar computation is exactly
  equation~\eqref{eq:app-marked-differential-datum}, including its signs and
  invariant marker tensors.

  In particular, a single fixed marker \(T\in\operatorname{End}(V){}_{\bar0}\)
  can occur in a full \(\lie{gl}(V)\)-equivariant marked differential only
  when \(T=\lambda\operatorname{id}_V\).  Its scalar shadow is then
  \[
    \hbar\,\lambda\,\operatorname{Str}(\operatorname{id}_V)\,
    \omega(f,g)\gamma_{\mathbf 1}=0 .
  \]
  Thus the nonzero \(P_+\) test of
  Proposition~\ref{prop:app-first-nonidentity-loop-scalar-symbol} is a
  genuine obstruction to a non-equivariant or externally marked
  enlargement, but it is not an allowed single-marker term in the full
  \(\lie{gl}(N|N)\)-equivariant Costello differential.  If one imposes
  only the even block-diagonal group, or supplies external marker
  tensors, the admissible datum is the invariant tensor family above, and
  the first nonzero scalar shadow is detected exactly by the displayed
  cyclic-supertrace coefficient.
\end{prop}

\begin{proof}
  The formula for
  \(\mathcal B^{(r)}_{\operatorname{st}=0}\) is
  Definition~\ref{def:app-supertraceless-monodromy-prehabitat} applied to
  the marked boundary expansion
  equation~\eqref{eq:app-marked-differential-datum}.  The scalar-shadow
  formula is the linear extension of
  Proposition~\ref{prop:app-first-nonidentity-loop-scalar-symbol}: a
  decomposable marker word contributes the cyclic supertrace of the
  ordered product around the open index loop, and an invariant tensor is a
  finite linear combination of such words.  Passing to cohomology replaces
  \(\omega\) by \([\bar c]\) as in
  Theorem~\ref{thm:app-supertraceless-monodromy-projection}.

  It remains to record the equivariance constraint.  A marked local
  functional with one fixed marker \(T\) transforms under \(g\in GL(V)\)
  by replacing the marker with \(gTg^{-1}\).  Equivariance of the
  differential therefore forces \(gTg^{-1}=T\), or infinitesimally
  \([A,T]=0\) for every \(A\in\lie{gl}(V)\).  Since \(T\) is even, write it
  in the parity decomposition as
  \(T=T_{\bar0}\oplus T_{\bar1}\).  Commutation with
  \(\lie{gl}(N)\oplus\lie{gl}(N)\) gives
  \(T_{\bar0}=\lambda_0\operatorname{id}\) and
  \(T_{\bar1}=\lambda_1\operatorname{id}\).  Commutation with the odd
  elementary endomorphisms forces \(\lambda_0=\lambda_1\).  Hence
  \(T=\lambda\operatorname{id}_V\).  The cyclic supertrace is then
  \(\lambda\operatorname{sdim}(V)=0\) for \(V=\C^{N|N}\).

  The final assertion is bookkeeping, not an additional vanishing
  theorem.  The only displayed \(d_{\mathrm{QME}}\) in the current
  Hamiltonian generator calculation is \(Q+\{I_0,-\}\) with no external
  endomorphism label, and
  Definition~\ref{def:app-balanced-costello-graph-subhabitat} requires
  identity-loop factorization after applying \(d_{\mathrm{QME}}\).  These
  terms have cyclic supertrace
  \(\operatorname{Str}(\operatorname{id}_V)=0\), so their
  \(\mathcal B_{\operatorname{st}=0}\)-class vanishes.  Any further
  marked graph/counterterm term is governed by its chosen oriented marked
  boundary complex; without that data,
  equation~\eqref{eq:app-marked-differential-datum} remains the residual.
\end{proof}

\begin{defn}[Finite-window marked Costello list]
\label{def:app-finite-window-marked-costello-list}
  Fix a finite-window graph/QME package for the local mixed HT
  brane-defect theory,
  \((I,w,M,\mathcal G_M)\) in the sense of
  Definition~\ref{defn:finite-window-graph-qme-package}.  A finite-window
  marked Costello list
  \[
    \mathcal G^{\mathrm{mk}}_M
  \]
  is a finite set of connected brane-defect graphs and local counterterm
  graphs with the following labels and boundary table.

  The admissible vertex labels are exactly the finite local labels visible
  in the chosen package:
  \begin{enumerate}[label=(L\arabic*)]
    \item external CE labels \(u_{a(t)dt\otimes f}\) and \(c_{B,\rho}\),
    with \(f,\rho\) in the Hamiltonian window \(M\), and with
    \(B\in\mathcal D^{\mathrm{reg}}_c(I)\) or
    \(B_\bullet\in\mathcal D'_{c,\mathrm{WF}}(I;\Gamma,M)\);
    \item interaction vertices from the finite-window local interaction
    \(I^w_{0,M}\);
    \item renormalized graph-weight vertices
    \(W^{R,M}_{\Gamma,L,w}(B^\Gamma_\bullet)\) and local counterterm
    vertices \(C^M_{\Gamma,w}(\epsilon;B^\Gamma_\bullet)\) supplied by
    the package;
    \item the central scalar-contact ghost \(\gamma_{\mathbf 1}\), only on
    the scalar-contact boundary vertices already used in
    Theorem~\ref{thm:app-scalar-contact-qme-class}.
  \end{enumerate}
  The admissible edge labels are the finite-scale local propagator
  \(P^{w,M}_{\epsilon,L}\), the corresponding regularized BV contraction
  \(\Delta^{w,M}_{\epsilon,L}\), and the boundary Weyl edge whose
  scalar part on two Hamiltonians is the formal-local scalar cocycle of
  Proposition~\ref{prop:app-local-geometry-scalar-shadow},
  \[
    (f,g)\longmapsto \omega(f,g)\gamma_{\mathbf 1}.
  \]
  The open index loops may carry marker tensors, defined below in
  Definition~\ref{def:app-oriented-full-equivariant-marked-boundary-complex}.

  For each \(\Gamma\in\mathcal G^{\mathrm{mk}}_M\) one specifies a finite
  ordered set
  \[
    \mathsf B_M(\Gamma)=\{\beta_1,\ldots,\beta_s\}
  \]
  of elementary codimension-one boundary operations.  If a differential,
  bracket, collision, or counterterm expansion is a finite sum, each
  summand is a separate \(\beta_i\).  Each operation has a coefficient
  \(\varepsilon_{\beta_i}\in\C[[\hbar]]\), an output graph
  \(\partial_{\beta_i}\Gamma\in\mathcal G^{\mathrm{mk}}_M\), and a
  marker-transport map \(\mu_{\beta_i}\).  The local-functional boundary
  operations are:
  \begin{enumerate}[label=(B\arabic*)]
    \item \emph{field differential:} if a vertex has label \(\ell_v\) and
    \(Q\ell_v=\sum_\alpha q_\alpha\ell_{v,\alpha}\), then the summand
    \((v,\alpha)\) replaces \(\ell_v\) by \(\ell_{v,\alpha}\) and has
    \(\varepsilon_{v,\alpha}=q_\alpha\), including the Koszul sign from
    moving \(Q\) to \(v\);
    \item \emph{BV edge contraction:} if an oriented internal edge \(e\)
    joins vertices \(v_1,v_2\) and
    \[
      \{\ell_{v_1},\ell_{v_2}\}_{\hbar,M}
        =
      \sum_\alpha q_{e,\alpha}\ell_{e,\alpha},
    \]
    then the summand \((e,\alpha)\) contracts \(e\), merges the two
    vertices with label \(\ell_{e,\alpha}\), and has coefficient
    \(q_{e,\alpha}\) multiplied by the ordinary BV orientation sign.  The
    bracket with \(I^w_{0,M}\) is this operation with one vertex labelled
    by \(I^w_{0,M}\);
    \item \emph{boundary collision/contact:} if an ordered boundary
    collision stratum \(S=(v_1,\ldots,v_k)\) has local expansion
    \[
      \operatorname{Coll}_{S}(\ell_{v_1},\ldots,\ell_{v_k})
        =
      \sum_\alpha q_{S,\alpha}\ell_{S,\alpha},
    \]
    then \((S,\alpha)\) collapses \(S\) to one boundary vertex labelled
    \(\ell_{S,\alpha}\).  For the two-Hamiltonian scalar contact this
    convention gives
    \[
      \operatorname{Coll}^{\mathrm{sc}}((a,f),(b,g))
        =
      \hbar\,\omega(f,g)\int_I a(t)b(t)\gamma_{\mathbf 1}(t)\,dt,
    \]
    antisymmetric in \((f,g)\);
    \item \emph{counterterm face:} if the small-diagonal extension of a
    graph subconfiguration has local counterterm expansion
    \(C^M_{\Gamma,w}(\epsilon;B_\bullet)=
    \sum_\alpha q_{\Gamma,\alpha}C^M_{\Gamma,\alpha}\), then
    \((\Gamma,\alpha)\) replaces that subconfiguration by the local
    counterterm vertex \(C^M_{\Gamma,\alpha}\);
  \end{enumerate}
  For the Hom-valued graph kernel one adjoins one further operation:
  the source differential face \(-\kappa d_{\mathrm{CE}}\), with the CE
  bracket coefficient included in \(\varepsilon_{\beta_i}\).  After
  evaluation on fixed CE cycles this extra face is absent.

  The finite list is required to contain every output of (B1)--(B4) that
  appears in the chosen package, and also the source face when the
  Hom-valued kernel table is used.  If an output graph, coefficient,
  marker transport, or counterterm vertex is not supplied in
  \(\mathcal G^{\mathrm{mk}}_M\), the missing datum is exactly that
  omitted boundary entry; no scalar-shadow theorem is asserted for the
  larger specialization.

  The orientation line used for signs is
  \[
    \operatorname{or}_M(\Gamma)
      =
    \det\mathsf B_M(\Gamma)\otimes
    \det E_{\mathrm{int}}(\Gamma)\otimes
    \bigotimes_v \operatorname{or}(\ell_v).
  \]
  The first determinant gives the incidence sign in the marked
  differential below.  The edge, vertex, BV, and CE Koszul signs are
  included in the coefficients \(\varepsilon_{\beta_i}\).

  The list is \emph{codimension-two closed} if, for every
  \(\beta_i\neq\beta_j\), the two composites are again entries of the
  finite list and, after the canonical identification of output orientation
  lines,
  \[
    \partial_{\beta_j|\beta_i}\partial_{\beta_i}\Gamma
      =
    \partial_{\beta_i|\beta_j}\partial_{\beta_j}\Gamma,
  \]
  \[
    \varepsilon_{\beta_i}\varepsilon_{\beta_j|\beta_i}
      =
    \varepsilon_{\beta_j}\varepsilon_{\beta_i|\beta_j},
    \qquad
    \mu_{\beta_j|\beta_i}\mu_{\beta_i}
      =
    \mu_{\beta_i|\beta_j}\mu_{\beta_j}.
    \label{eq:app-codim-two-marked-closure}
  \]
  Here \(\beta_j|\beta_i\) denotes the surviving operation induced by
  \(\beta_j\) after applying \(\beta_i\).  Thus, from a seed graph
  \(\Gamma\), the actual finite window consists of the finitely many
  faces obtained by deleting subsets of the finite set
  \(\mathsf B_M(\Gamma)\):
  \[
    \mathcal G^{\mathrm{mk}}_M
      =
    \left\{
      \partial_J\Gamma\ \middle|
      \Gamma\in\mathcal S_M,\ J\subset\mathsf B_M(\Gamma)
    \right\},
    \qquad
    \partial_J=\partial_{\beta_{j_k}}\cdots\partial_{\beta_{j_1}} .
  \]
  Here \(\mathcal S_M\) is the finite seed set of package graphs and
  counterterm graphs, and the codimension-two closure equations make
  \(\partial_J\) independent of the chosen order of \(J\).  Failure of this
  displayed equality is precisely failure to supply the finite graph list.
\end{defn}

\begin{defn}[Oriented full-equivariant marked boundary complex]
\label{def:app-oriented-full-equivariant-marked-boundary-complex}
  Let \(V=\C^{N|N}\).  For \(m\geq1\), let
  \[
    \mathcal M_m^{\mathrm{full}}(V)
      \subset
    \bigl(\operatorname{End}(V)^{\otimes m}\bigr)^{\lie{gl}(V)}
  \]
  be the finite span of the signed Brauer permutation tensors
  \(\mathfrak C_\sigma\), \(\sigma\in S_m\).  Here
  \(\mathfrak C_\sigma\) is the string-diagram tensor obtained by
  writing
  \(\operatorname{End}(V)=V\otimes V^\vee\), pairing the \(V^\vee\)
  factor in the \(i\)-th copy with the \(V\) factor in the
  \(\sigma(i)\)-th copy, and using the Koszul signs of the symmetric
  monoidal category of super vector spaces.  For \(m=0\) the marker is
  the empty word.

  A full-equivariant marked graph/counterterm generator in a finite
  window \(M\) is a tuple
  \[
    X_{\Gamma,o,\mathfrak m,M},
  \]
  where \(\Gamma\in\mathcal G^{\mathrm{mk}}_M\) for a finite-window
  marked Costello list of
  Definition~\ref{def:app-finite-window-marked-costello-list},
  \(o\) is an orientation of the finite set
  \(\mathsf B_M(\Gamma)\), and
  \(\mathfrak m=(\mathfrak m_L)_L\) assigns to each marked scalar open
  index loop \(L\) a tensor
  \(\mathfrak m_L\in\mathcal M_{m_L}^{\mathrm{full}}(V)\).

  If
  \[
    o=\beta_1\wedge\cdots\wedge\beta_s
  \]
  is the ordered orientation of \(\mathsf B_M(\Gamma)\), put
  \(o/\beta_i=\beta_1\wedge\cdots\widehat{\beta_i}\cdots\wedge\beta_s\).
  The marker-transport map
  \[
    \mu_{\beta_i}\colon
    \bigotimes_L\mathcal M_{m_L}^{\mathrm{full}}(V)
    \longrightarrow
    \bigotimes_{L'}\mathcal M_{m'_{L'}}^{\mathrm{full}}(V)
  \]
  is obtained by the signed evaluation, coevaluation, braiding, loop
  splitting, and loop concatenation maps that define the open index
  contractions.  Concretely, loop merging sends two cyclic words to their
  boundary-ordered concatenation with the Koszul braiding sign; loop
  splitting inserts the coevaluation tensor at the split points; deletion
  of a closed marked loop applies the cyclic supertrace.  The signed
  marked differential is
  \[
    d_{\mathrm{mk}}X_{\Gamma,o,\mathfrak m,M}
      =
    \sum_{i=1}^{s}
      (-1)^{i-1}
      \varepsilon_{\beta_i}\,
      X_{\partial_{\beta_i}\Gamma,\;o/\beta_i,\;
        \mu_{\beta_i}(\mathfrak m),\;M}.
    \label{eq:app-oriented-marked-differential}
  \]
  Its scalar-contact filtration is
  \[
    p(\Gamma,\mathfrak m)
      =
    \#\{\text{scalar-contact open loops in }\Gamma\},
    \qquad
    F^p=\overline{\operatorname{Span}}\{X_{\Gamma,o,\mathfrak m,M}
      :p(\Gamma,\mathfrak m)\geq p\}.
  \]
  For an elementary boundary operation set
  \[
    r(\beta_i)
      =
    p(\Gamma,\mathfrak m)
      -p(\partial_{\beta_i}\Gamma,\mu_{\beta_i}(\mathfrak m)).
  \]
  The first filtration-lowering marked boundary is therefore
  \[
    \left[d_{\mathrm{mk}}X_{\Gamma,o,\mathfrak m,M}
    \right]_{F^{p-r}/F^{p-r+1}}
      =
    \sum_{i:\,r(\beta_i)=r}
      (-1)^{i-1}
      \varepsilon_{\beta_i}\,
      X_{\partial_{\beta_i}\Gamma,\;o/\beta_i,\;
        \mu_{\beta_i}(\mathfrak m),\;M}.
    \label{eq:app-oriented-marked-first-lowering}
  \]
\end{defn}

\begin{lem}[Incidence signs]\label{lem:app-oriented-marked-incidence-signs}
  If the finite marked Costello list is codimension-two closed in the sense
  of Definition~\ref{def:app-finite-window-marked-costello-list}, then
  \(d_{\mathrm{mk}}^2=0\).
\end{lem}

\begin{proof}
  This is the cellular boundary sign computation.  The codimension-two
  output obtained from \(\beta_i\) and \(\beta_j\), \(i<j\), appears
  twice.  The order \(\beta_i\) then \(\beta_j\) contributes
  \[
    (-1)^{i-1}(-1)^{j-2}
    \varepsilon_{\beta_i}\varepsilon_{\beta_j|\beta_i},
  \]
  because \(\beta_i\) has been removed before \(\beta_j\) is counted.  The
  order \(\beta_j\) then \(\beta_i\) contributes
  \[
    (-1)^{j-1}(-1)^{i-1}
    \varepsilon_{\beta_j}\varepsilon_{\beta_i|\beta_j}.
  \]
  The incidence signs differ by a minus sign, while
  \eqref{eq:app-codim-two-marked-closure} identifies the coefficients,
  output graph, and marker tensor.  Hence the two terms cancel.  Summing
  over all unordered pairs gives \(d_{\mathrm{mk}}^2=0\).
\end{proof}

\begin{prop}[Finite-window marked habitat compatibility]
\label{prop:app-finite-window-marked-habitat-compatibility}
  Let \((I,w,M,\mathcal G_M)\) be a finite-window graph/QME package for
  the local mixed HT brane-defect theory and let
  \(\mathcal G^{\mathrm{mk}}_M\) be a codimension-two closed finite-window
  marked Costello list whose labels are the admissible local labels of
  Definition~\ref{def:app-finite-window-marked-costello-list}.  Let
  \[
    \mathfrak G^\bullet_{\mathrm{mk,full},M}(I)
  \]
  be the completed span of the full-equivariant marked generators
  \(X_{\Gamma,o,\mathfrak m,M}\).  Then
  \[
    \mathfrak G^\bullet_{\mathrm{mk,full},M}(I)
      \subset
    \mathfrak Q^\bullet_{\mathcal G,M}(I)
  \]
  and \(d_{\mathrm{mk}}\) is the restriction of the finite-window
  differential \(d_M=Q+\{I^w_{0,M},-\}_{\hbar,M}\) on the local-functional
  table.  On the Hom-valued kernel table, adjoining the source face makes
  the same formula the convolution differential
  \(d_{\mathbf K}=d_M-(-1)^{|\cdot|}(\cdot)d_{\mathrm{CE}}\) before
  evaluation.  The graph kernel
  \(\kappa^M_{\mathcal G,w,I}\) and the local counterterms of the package
  land in this marked habitat exactly when their graph components and
  counterterm faces occur in \(\mathcal G^{\mathrm{mk}}_M\).  Truncation
  \(\pi_{M,N}\) and admissible weight transport \(R^M_{w,w'}\) preserve the
  marked habitat precisely when they send the finite boundary table,
  coefficients, and marker transports to the corresponding lower-window or
  transported table.
\end{prop}

\begin{proof}
  The vertex and edge labels are those of
  Definition~\ref{defn:finite-window-graph-qme-package}: regular-density
  or wavefront-admissible graph labels, finite-window propagators,
  renormalized weights, and local counterterms.  Hence the represented
  local functionals lie in
  \(\mathfrak Q^\bullet_{\mathcal G,M}(I)\).  The boundary table
  (B1)--(B4) is exactly the expansion of \(Q\), the BV bracket with
  \(I^w_{0,M}\), collision/contact faces, and Costello counterterm faces on
  these generators.  Thus the signed sum
  \eqref{eq:app-oriented-marked-differential} is \(d_M\) on the local
  functional span.  In the convolution kernel before CE evaluation, the
  additional source face is the missing source differential term in
  \(d_{\mathbf K}=d_M-(-1)^{|\cdot|}(\cdot)d_{\mathrm{CE}}\); for the
  degree-zero kernel this is \(-\kappa d_{\mathrm{CE}}\).

  The inclusion of the graph kernel and counterterms is tautological from
  the same finite label list.  If a graph component or counterterm face is
  absent, the equality with \(d_M\) fails at that named face, which is the
  exact missing datum.  The finite-window package requires compatibility of
  graph weights, counterterms, scalar projections, and kernels with
  \(\pi_{M,N}\) and \(R^M_{w,w'}\); applying this requirement to each entry
  of the finite boundary table gives the stated preservation criterion.
\end{proof}

\begin{lem}[Cyclic supertrace of full-equivariant marker tensors]
\label{lem:app-full-equivariant-cyclic-supertrace}
  Let \(\tau_m=(1\,2\,\cdots\,m)\).  For the signed Brauer permutation
  tensor \(\mathfrak C_\sigma\in\mathcal M_m^{\mathrm{full}}(V)\),
  \[
    \operatorname{Str}_{\mathrm{cyc}}(\mathfrak C_\sigma)
      =
    \bigl(\operatorname{sdim}V\bigr)^{c(\tau_m^{-1}\sigma)},
    \label{eq:app-strcyc-brauer-permutation}
  \]
  where \(c(\pi)\) is the number of cycles of the permutation \(\pi\).
  Hence, for \(V=\C^{N|N}\),
  \[
    \operatorname{Str}_{\mathrm{cyc}}(\mathfrak m)=0
    \qquad
    \text{for every }
    \mathfrak m\in\mathcal M_m^{\mathrm{full}}(V)
  \]
  and also for the empty marker word.
\end{lem}

\begin{proof}
  Compose the string diagram for \(\mathfrak C_\sigma\) with the cyclic
  supertrace functional
  \((T_1,\ldots,T_m)\mapsto\operatorname{Str}(T_1\cdots T_m)\).  The
  cyclic multiplication identifies the output of the \(i\)-th
  endomorphism with the input of the \(\tau_m(i)\)-th endomorphism,
  while \(\mathfrak C_\sigma\) identifies it with the input of the
  \(\sigma(i)\)-th endomorphism.  The resulting closed diagram has
  \(c(\tau_m^{-1}\sigma)\) closed supertraces.  Each closed supertrace
  evaluates to
  \(\operatorname{Str}(\operatorname{id}_V)=\operatorname{sdim}V\).
  The signs in \(\mathfrak C_\sigma\) are exactly the Koszul signs needed
  for this diagrammatic evaluation, giving
  \eqref{eq:app-strcyc-brauer-permutation}.  Since
  \(\operatorname{sdim}\C^{N|N}=0\), every nonempty full-equivariant
  marker tensor has zero cyclic supertrace.  The empty word gives the
  identity-loop coefficient
  \(\operatorname{Str}(\operatorname{id}_V)=0\).
\end{proof}

\begin{thm}[Full-equivariant marked scalar shadow vanishing]
\label{thm:app-full-equivariant-marked-shadow-vanishing}
  Let \(V=\C^{N|N}\).  Let
  \(\mathfrak G^\bullet_{\mathrm{mk,full}}(I)\) be a complete
  finite-window inverse-limit graph/counterterm habitat for the
  formal-local mixed HT brane-defect theory generated
  windowwise by codimension-two closed finite-window marked Costello lists
  of Definition~\ref{def:app-finite-window-marked-costello-list} and by the
  full-equivariant marked generators of
  Definition~\ref{def:app-oriented-full-equivariant-marked-boundary-complex},
  compatible with truncation and admissible weight transport as in
  Proposition~\ref{prop:app-finite-window-marked-habitat-compatibility}.
  This is a finite-list statement in every window; no summation over all
  graph topologies, compact target, or global BCOV datum is part of the
  hypothesis.  Then the marked preservation obstruction has zero
  cyclic-supertrace shadow:
  \[
    \operatorname{Str}_{\mathrm{cyc}}
      \bigl(\mu_{\beta_i}(\mathfrak m)\bigr)=0
    \qquad\text{for every elementary boundary term } \beta_i.
  \]
  Consequently, for every homogeneous
  \(X_{\Gamma,o,\mathfrak m,M}\in F^p\), the first scalar CE shadow is
  \[
    \hbar
    \sum_{i:\,r(\beta_i)=r}
      (-1)^{i-1}
      \varepsilon_{\beta_i}\,
      \operatorname{Str}_{\mathrm{cyc}}
        \bigl(\mu_{\beta_i}(\mathfrak m)\bigr)
      \omega(f_{\partial_{\beta_i}\Gamma},
             g_{\partial_{\beta_i}\Gamma})
      \gamma_{\mathbf 1}
      =0,
  \]
  and its cohomology class on the \([\bar c]\)-line is
  \[
    \hbar
    \sum_{i:\,r(\beta_i)=r}
      (-1)^{i-1}
      \varepsilon_{\beta_i}\,
      \operatorname{Str}_{\mathrm{cyc}}
        \bigl(\mu_{\beta_i}(\mathfrak m)\bigr)
      [\bar c]\,\gamma_{\mathbf 1}
      =0 .
  \]
  Thus the zero scalar projection is a filtered chain map on
  \(\mathfrak G^\bullet_{\mathrm{mk,full}}(I)\), and all scalar
  projection obstruction classes vanish on this full-equivariant marked
  graph/counterterm habitat.
\end{thm}

\begin{proof}
  The marker transport maps \(\mu_{\beta_i}\) are built from the
  evaluation, coevaluation, braiding, loop splitting, and loop
  concatenation maps of the \(\lie{gl}(V)\)-module \(V\).  Therefore they
  send full-equivariant signed Brauer marker tensors to full-equivariant
  signed Brauer marker tensors.  Lemma~\ref{lem:app-full-equivariant-cyclic-supertrace}
  gives zero cyclic supertrace for every such output tensor when
  \(V=\C^{N|N}\).

  Formula~\eqref{eq:app-oriented-marked-first-lowering} gives the actual
  signed first filtration-lowering boundary.  Applying the one-loop
  scalar computation of
  Proposition~\ref{prop:app-first-nonidentity-loop-scalar-symbol} to each
  boundary term gives exactly the displayed scalar CE representative.
  Every coefficient in that sum is zero by the preceding paragraph.
  Passing to cohomology replaces \(\omega\) by the nonzero class
  \([\bar c]\), but it is multiplied by the same zero coefficient.  Thus
  \(\widehat\sigma=0\) satisfies
  \(d_{\mathrm{CE}}\widehat\sigma-\widehat\sigma d_{\mathrm{mk}}=0\).
  Proposition~\ref{prop:app-finite-window-marked-habitat-compatibility}
  identifies \(d_{\mathrm{mk}}\) with the restricted finite-window QME
  differential on the chosen finite table.  Hence this is the required
  filtered chain projection, and
  Lemma~\ref{lem:filtered-scalar-projection-obstruction} kills the scalar
  projection tower on the habitat.
\end{proof}

\begin{defn}[Non-scalar kernel complex for a marked finite window]
\label{def:app-nonscalar-kernel-row-complex}
  Fix a finite-window graph/QME package
  \((I,w,M,\mathcal G_M)\), a codimension-two closed marked Costello list
  \(\mathcal G^{\mathrm{mk}}_M\), and a scalar projection
  \[
    \widehat\sigma_{\mathrm{sc},M}\colon
    \mathfrak Q^\bullet_{\mathcal G,M}(I)
      \longrightarrow
    S^\bullet_M,\qquad
    S^\bullet_M=
    C^\bullet_{\mathrm{Lie}}(\bar A;\C\gamma_{\mathbf 1})[1][[\hbar]] .
  \]
  The non-scalar obstruction complex in that window is
  \[
    \mathcal K^\bullet_{\mathrm{ns},M}(I)
      =
    \ker\widehat\sigma_{\mathrm{sc},M},
    \qquad
    d_{\mathrm{ns},M}=d_M|_{\mathcal K_{\mathrm{ns},M}} .
  \]
  It is a complex precisely when \(\widehat\sigma_{\mathrm{sc},M}\) is a
  filtered chain map.  In the full-equivariant marked habitat of
  Theorem~\ref{thm:app-full-equivariant-marked-shadow-vanishing} one has
  \(\widehat\sigma_{\mathrm{sc},M}=0\), hence
  \[
    \mathcal K^\bullet_{\mathrm{ns},M}(I)
      =
    \mathfrak G^\bullet_{\mathrm{mk,full},M}(I).
  \]

  Let \(\mathsf A^{\mathrm{fw}}_{n,M}\) be a valid row array in the sense
  of Definition~\ref{defn:finite-window-graph-array}.  If the scalar gate
  vanishes,
  \[
    \sum_{\alpha\in\mathsf A^{\mathrm{fw}}_{n,M}}S_{\alpha,M}=0,
  \]
  the order-\(n\) kernel residual is the scalar-zero row sum
  \[
    R^{\mathrm{ns}}_{n,M}
      =
    \sum_{\alpha\in\mathsf A^{\mathrm{fw}}_{n,M}}
      R^{\mathrm{row}}_{\alpha,M}
      \in\mathcal K^1_{\mathrm{ns},M}(I).
  \]
  Its obstruction class is
  \[
    \mathfrak o^{\mathrm{ns}}_{n,M}
      =
    [R^{\mathrm{ns}}_{n,M}]
      \in
    H^1\!\left(
      \mathcal K^\bullet_{\mathrm{ns},M}(I),d_{\mathrm{ns},M}
    \right).
  \]
  A degree-zero counterterm
  \(C_{n,M}\in\mathcal K^0_{\mathrm{ns},M}(I)\) exists at this order if
  and only if
  \[
    d_M C_{n,M}=-R^{\mathrm{ns}}_{n,M},
    \qquad
    \mathfrak o^{\mathrm{ns}}_{n,M}=0.
  \]
\end{defn}

\begin{defn}[Native finite-window realization habitat]
\label{def:app-native-finite-window-realization-habitat}
  A native Costello/QME realization habitat for the mixed HT brane-defect
  model on
  \[
    \R^2_{\mathrm{top}}\times\widehat{\C^2}_{0,\mathrm{hol}}
  \]
  is the following finite-window datum.  First, one fixes a graph/QME
  package \((I,w,M,\mathcal G_M)\) in the sense of
  Definition~\ref{defn:finite-window-graph-qme-package}, a
  codimension-two closed marked list
  \(\mathcal G^{\mathrm{mk}}_M\) in the sense of
  Definition~\ref{def:app-finite-window-marked-costello-list}, and the
  completed local-functional complex
  \[
    \mathfrak Q^\bullet_{\mathcal G,M}(I)
      =
    \left(
      \mathcal O^{\mathrm{bd}}_{\mathrm{loc},\mathcal G,M}
      (\mathcal E_w^M|_I;
       A^{\mathrm{pp},w,M}_{\partial,\hbar}(I)),
      d_M
    \right),
    \qquad
    d_M=Q+\{I^w_{0,M},-\}_{\hbar,M}.
  \]
  The topology is the complete topology generated by the \(\hbar\)-adic,
  scalar-contact, and finite-window filtrations.  The complex is part of
  the datum: every \(d_M\)-boundary, collision face, BV bracket row, and
  local counterterm face needed for closure must occur in the finite
  table.

  Second, the source is the labelled finite CE object
  \(C^\bullet_{\mathrm{lab}}(\mathcal G_M)\).  The Hom-valued kernel
  complex is
  \[
    \mathbf K^\bullet_{\mathcal G,M}(I)
      =
    \operatorname{Hom}_{\mathrm{cont}}
    \left(
      C^\bullet_{\mathrm{lab}}(\mathcal G_M),
      \mathfrak Q^\bullet_{\mathcal G,M}(I)
    \right),
  \]
  with
  \[
    d_{\mathbf K,M}T
      =
    d_M T-(-1)^{|T|}T\,d_{\mathrm{CE}},
    \qquad
    \operatorname{Curv}_{\mathbf K,M}(\kappa)
      =
    d_{\mathbf K,M}\kappa+\frac12[\kappa,\kappa]_{\mathbf K,M}.
  \]
  Thus the source-face term is not optional.  For a degree-zero kernel it
  contributes \(-\kappa d_{\mathrm{CE}}\).

  Third, the kernel has a fixed zero-edge part:
  \[
    \kappa^{\mathrm{coef},M}_{\hbar,w,I}(u_{a(t)dt\otimes f})
      =
    \iota_I(B^{w,M}_f(a)),
    \qquad
    \kappa^{\mathrm{coef},M}_{\hbar,w,I}(c_{B,\rho})
      =
    \iota_I(\Theta^{w,M}_\rho(B)).
  \]
  Here \(a\in C^\infty_c(I)\), \(B\) is a regular compactly supported
  density or a graphwise wavefront-admissible current for the chosen
  finite row, and \(f,\rho\) lie in the window \(M\).  Positive-edge
  components are exactly the renormalized graph weights
  \(K^M_\Gamma=W^{R,M}_{\Gamma,L,w}(B^\Gamma_\bullet)\) supplied by the
  package.  No graph theorem for arbitrary \(B\in\mathcal D'_c(I)\) is
  included in this definition.

  Fourth, one supplies a filtered chain scalar projection
  \[
    \widehat\sigma_{\mathrm{sc},M}\colon
    \mathfrak Q^\bullet_{\mathcal G,M}(I)
      \longrightarrow
    C^\bullet_{\mathrm{Lie}}(\bar A;\C\gamma_{\mathbf 1})[1][[\hbar]],
  \]
  and puts
  \[
    \mathcal K^\bullet_{\mathrm{ns},M}(I)
      =
    \ker\widehat\sigma_{\mathrm{sc},M}.
  \]
  The ambient QME complex is always denoted
  \(\mathfrak Q^\bullet_{\mathcal G,M}(I)\).  If the windowed
  scalar-zero QME notation
  \(\mathcal Q^\bullet_{w,M}(I)\) is used for the same kernel, then
  \[
    \mathcal Q^\bullet_{w,M}(I)
      =
    \mathcal K^\bullet_{\mathrm{ns},M}(I)
      =
    \ker\widehat\sigma_{\mathrm{sc},M}.
  \]
  In this appendix \(\mathcal K_{\mathrm{ns},M}\) is used for the
  non-scalar kernel, so that it is not confused with the ambient QME
  complex \(\mathfrak Q_{\mathcal G,M}\).
  If this projection is not a chain map, the non-scalar kernel complex is
  not formed.
\end{defn}

\begin{thm}[Constructive non-scalar Costello/QME realization criterion]
\label{thm:app-constructive-nonscalar-costello-qme-realization}
  Work in a native finite-window realization habitat of
  Definition~\ref{def:app-native-finite-window-realization-habitat}.
  Suppose that the QME has been solved through order \(\hbar^{n-1}\) by
  counterterms \(C_{<n,M}\).  The order-\(n\) realization problem is
  decided by the following finite row construction.

  For every valid row array
  \(\mathsf A^{\mathrm{fw}}_{n,M}\) of
  Definition~\ref{defn:finite-window-graph-array}, the rows are precisely
  \[
    R^{\mathrm{row}}_{d\Gamma,M}
      =
    \varepsilon_\Gamma\,d_M K^M_\Gamma,
  \]
  \[
    R^{\mathrm{row}}_{\mathrm{CE}(\Gamma,\nu),M}
      =
    -\varepsilon^{\mathrm{CE}}_{\Gamma,\nu}
      K^M_\Gamma(\zeta_{M,\nu}),
  \]
  \[
    R^{\mathrm{row}}_{b(\Gamma_1,\Gamma_2),M}
      =
    \frac12\varepsilon_{\Gamma_1,\Gamma_2}
      \{K^M_{\Gamma_1},K^M_{\Gamma_2}\}_{\hbar,M},
  \]
  together with the homogeneous lower-counterterm rows in
  \[
    [\hbar^n]\left(
      d_M C_{<n,M}
      +\frac12\{C_{<n,M},C_{<n,M}\}_{\hbar,M}
    \right).
  \]
  Their scalar gates are
  \[
    S_{\alpha,M}
      =
    \widehat\sigma_{\mathrm{sc},M}
      (R^{\mathrm{row}}_{\alpha,M}),
  \]
  and the scalar gate for the order is the single equation
  \[
    \mathfrak s_{n,M}
      =
    \sum_{\alpha\in\mathsf A^{\mathrm{fw}}_{n,M}}S_{\alpha,M}=0.
    \label{eq:app-native-scalar-gate}
  \]
  If \eqref{eq:app-native-scalar-gate} fails, the non-scalar problem has
  not started at order \(n\).  If it holds, then
  \[
    R^{\mathrm{ns}}_{n,M}
      =
    \sum_{\alpha\in\mathsf A^{\mathrm{fw}}_{n,M}}
      R^{\mathrm{row}}_{\alpha,M}
      \in \mathcal K^1_{\mathrm{ns},M}(I)
  \]
  is a cocycle whenever the finite row table is codimension-two closed and
  the lower orders have been solved.  Its obstruction class is
  \[
    \mathfrak o^{\mathrm{ns}}_{n,M}
      =
    [R^{\mathrm{ns}}_{n,M}]
      \in
    H^1(\mathcal K^\bullet_{\mathrm{ns},M}(I),d_{\mathrm{ns},M}).
  \]

  Choose degree-one row coordinates
  \(\rho^M_1,\ldots,\rho^M_a\) and degree-zero counterterm candidates
  \(\eta^M_1,\ldots,\eta^M_b\) inside the same finite package, and write
  \[
    R^{\mathrm{ns}}_{n,M}=\sum_i r^M_i\rho^M_i,
    \qquad
    d_{\mathrm{ns},M}\eta^M_j=\sum_i A^M_{ij}\rho^M_i .
  \]
  A fixed-window counterterm in the supplied span exists exactly when
  \[
    A^M c=-r^M
    \label{eq:app-native-counterterm-linear-system}
  \]
  is solvable.  If it is not solvable, any covector
  \[
    \ell A^M=0,\qquad \ell(r^M)\neq0
  \]
  is a checkable certificate that
  \(\mathfrak o^{\mathrm{ns}}_{n,M}\neq0\) in the displayed finite row
  subcomplex.  A compatible finite-window tower also requires truncation
  matrices \(T^{M,N}\) on the rows and \(P^{M,N}\) on candidate primitives
  satisfying
  \[
    T^{M,N}r^M=r^N,\qquad
    T^{M,N}A^M=A^N P^{M,N},
  \]
  and the Roos class
  \[
    [P^{M,N}c_M-c_N]\in
    H^0(\mathcal K^\bullet_{\mathrm{ns},N}(I))
  \]
  must vanish.

  Iterating these finite-window tests for all \(n\), with graphwise
  regular-density or wavefront-admissible kernel data and with compatible
  counterterms, is equivalent to
  \[
    \operatorname{Curv}_{\mathbf K}
      (\kappa_{\mathcal G,w,I}+C_{\hbar,w})=0
  \]
  in the completed native brane-defect kernel complex.  This is the
  strongest realization statement supplied by the current manuscript.  A
  larger Costello specialization must add the missing graph weight, source
  face, scalar gate, primitive candidate, or truncation matrix; none is
  supplied by analogy with compact BCOV theory.
\end{thm}

\begin{proof}
  The first three displayed row formulas are the homogeneous terms in
  \[
    d_M\kappa^M_{\mathcal G,w,I}
      -\kappa^M_{\mathcal G,w,I}d_{\mathrm{CE}}
      +\frac12[\kappa^M_{\mathcal G,w,I},
               \kappa^M_{\mathcal G,w,I}]_{\mathbf K,M}.
  \]
  The minus sign in the source-face row is the term
  \(-\kappa d_{\mathrm{CE}}\) in
  Definition~\ref{def:app-native-finite-window-realization-habitat}.  The
  last displayed family is the order-\(n\) contribution of lower
  counterterms to the same curved Maurer--Cartan equation.

  Applying the chain scalar projection gives the scalar gate.  If the
  gate is nonzero, the residual is not in
  \(\ker\widehat\sigma_{\mathrm{sc},M}\).  If the gate is zero, the
  residual is in the kernel complex.  Codimension-two closure is exactly
  the finite-row form of the Bianchi identity: the two ways of taking two
  elementary faces have opposite incidence signs and the same output graph,
  coefficient, and marker transport.  Since lower orders are solved, this
  makes \(R^{\mathrm{ns}}_{n,M}\) a \(d_{\mathrm{ns},M}\)-cocycle.

  A counterterm \(C_{n,M}\) changes the residual by
  \(d_{\mathrm{ns},M}C_{n,M}\).  Expanding \(C_{n,M}\) in the supplied
  degree-zero row basis gives the linear system
  \eqref{eq:app-native-counterterm-linear-system}.  Failure of solvability
  is equivalent to the residual vector being outside the image of \(A^M\);
  a left-null covector with nonzero value on \(r^M\) is precisely the
  finite-dimensional cokernel certificate.  The two truncation identities
  say that both residuals and boundaries commute with finite-window
  projection.  After windowwise primitives are chosen, their failure to
  commute with truncation is the displayed Roos \(1\)-cocycle in
  \(H^0\).  Thus the fixed-window \(H^1\) class and the Roos primitive
  class are exactly the obstruction pair.  Passing to all orders gives the
  completed curvature equation.
\end{proof}

\begin{prop}[Closed rows and the first supplied open source-face row]
\label{prop:app-closed-rows-and-theta-three-source-face}
  In the full-equivariant branch \(V=\C^{N|N}\), the one-cross
  two-Hamiltonian scalar-contact row
  \(\alpha_{\mathrm{sc}}\) with inputs \((z_1,z_2)\) has
  \[
    R^{\mathrm{row}}_{\alpha_{\mathrm{sc}},M}=0,
    \qquad
    S_{\alpha_{\mathrm{sc}},M}=0,
    \qquad
    C_{\alpha_{\mathrm{sc}},M}=0 .
  \]
  The labelled scalar-zero row \(\alpha_{11}\) with inputs
  \((z_1,z_1)\) has
  \[
    R^{\mathrm{row}}_{\alpha_{11},M}
      =
    \operatorname{Str}_{\mathrm{cyc}}
      (\mu_{\beta_{11}}(\mathfrak m))\,
    \omega(z_1,z_1)
    \int_Ia(t)b(t)\gamma_{\mathbf 1}(t)\,dt
      =0,
  \]
  \[
    S_{\alpha_{11},M}=0,\qquad C_{\alpha_{11},M}=0,
  \]
  with truncation \(t^{M,N}_{\alpha_{11}\alpha_{11}}=1\) when the labelled
  zero row is retained and the lower window still contains \(z_1\).

  The first supplied nonzero scalar-zero source-face row in the present
  row machinery is the order-three symbolic row
  \(\theta_3=\mathrm{CE}(\Theta_3,\nu_3)\), with
  \[
    R^{\mathrm{row}}_{\theta_3,M}
      =
    -K^M_{\Theta_3}(\zeta_{M,\nu_3})
      =:\mathsf E_{\theta_3,M},
    \qquad
    S_{\theta_3,M}=0 .
  \]
  If the row-closure datum
  \(d_{\mathrm{ns},M}\mathsf E_{\theta_3,M}=0\) is supplied, it defines
  \[
    \mathfrak o^{\mathrm{ns}}_{\theta_3,M}
      =
    [\mathsf E_{\theta_3,M}]
    \in
    H^1(\mathcal K^\bullet_{\mathrm{ns},M}(I),d_{\mathrm{ns},M}).
  \]
  In the finite row subcomplex currently supplied for this row,
  \[
    \mathcal K^0_{\theta_3,M}=0,\qquad
    \mathcal K^1_{\theta_3,M}=\C\mathsf E_{\theta_3,M},\qquad
    d_{\theta_3,M}=0,
  \]
  the covector
  \[
    \ell_{\theta_3}(\mathsf E_{\theta_3,M})=1
  \]
  certifies a nonzero finite-row obstruction.  The exact missing healing
  datum is either a degree-zero local counterterm
  \(C_{\theta_3,M}\) with
  \[
    d_{\mathrm{ns},M}C_{\theta_3,M}
      =-\mathsf E_{\theta_3,M},
  \]
  compatible under finite-window truncation, or additional Hom-valued
  companion rows whose sum changes the residual vector.  Without one of
  these data the order-three larger package is not a solved QME row.
\end{prop}

\begin{proof}
  The row \(\alpha_{\mathrm{sc}}\) is zero as a local functional by the
  cyclic supertrace computation of
  Lemma~\ref{lem:app-full-equivariant-cyclic-supertrace}; this is the
  closed row of Proposition~\ref{prop:first-finite-window-array-rows}.
  The row \(\alpha_{11}\) is zero for the independent reason
  \[
    \omega(z_1,z_1)=
    [\{z_1,z_1\}_{\Omega_{\mathrm{loc}}}]_0=0,
  \]
  which is Proposition~\ref{prop:app-first-scalar-zero-order-one-row}.
  In both cases the zero counterterm is a primitive and the stated
  truncation identities send zero rows to zero rows.

  The formula for \(\theta_3\) is the source-face row
  \(-\kappa d_{\mathrm{CE}}\) from
  Theorem~\ref{thm:app-constructive-nonscalar-costello-qme-realization},
  with the concrete row labels of
  Proposition~\ref{prop:concrete-order-three-larger-package-row}.  Its
  scalar gate is zero in the full-equivariant marked habitat.  In the
  one-row finite complex, there is no degree-zero candidate primitive, so
  the boundary matrix is the \(1\times0\) zero matrix and the residual
  vector is the basis vector \(\mathsf E_{\theta_3,M}\).  The displayed
  covector annihilates the image of the differential and not the residual.
  This proves nonvanishing only in the supplied finite row subcomplex.
  The displayed \(C_{\theta_3,M}\) or companion-row datum is exactly what
  would enlarge the complex enough to heal the row.
\end{proof}

\begin{prop}[Bianchi-exposed lower gate for \(\theta_3\)]
\label{prop:app-theta-three-bianchi-exposed-lower-gate}
  In the lower window \(N=2\), let
  \(E^2_{\nu_{02}}\) and \(E^2_{\nu_{20}}\) be the two transported
  source-face rows attached to the \(\theta_3\) source differential.  The
  canonical source-coordinate candidate is
  \[
    \zeta^0_2=u_{a,H_{1,0}}u_{b,H_{0,1}},
    \qquad
    D^2_{02,20}=K^2_{\Theta_3}(\zeta^0_2),
  \]
  with
  \[
    d_{\mathrm{CE}}\zeta^0_2
      =
    2\zeta_{\nu_{02}}-2\zeta_{\nu_{20}} .
  \]
  The Hom-valued Bianchi defect is
  \[
    \mathfrak b^2_{02,20}
      =
    d_{\mathrm{ns},2}K^2_{\Theta_3}(\zeta^0_2)
      -
    K^2_{\Theta_3}(d_{\mathrm{CE}}\zeta^0_2).
  \]
  In the Bianchi-exposed basis
  \[
    (E^2_{\nu_{02}},E^2_{\nu_{20}},\mathfrak b^2_{02,20}),
  \]
  the current source-coordinate column and the lower transported residual
  are
  \[
    A_D^2=(-2,2,1)^T,\qquad r_2=(2,-2,0)^T .
  \]
  Hence
  \[
    A_D^2c=-r_2
  \]
  has no solution in the supplied lower local-functional row habitat.  The
  normalized cokernel
  \[
    \widetilde\lambda_{02,\mathfrak b}
      =
    \frac12(E^2_{\nu_{02}})^*
      +(\mathfrak b^2_{02,20})^*
  \]
  satisfies
  \[
    \widetilde\lambda_{02,\mathfrak b}A_D^2=0,\qquad
    \widetilde\lambda_{02,\mathfrak b}(r_2)=1 .
  \]

  The missing Bianchi-counterterm target is
  \[
    -\mathfrak b^2_{02,20}=(0,0,-1)^T,
  \]
  and it is not in the image of the current column \(A_D^2\), since
  \[
    \widetilde\lambda_{02,\mathfrak b}((0,0,-1)^T)=-1 .
  \]
  A repair must therefore add a genuine scalar-zero local-functional
  column
  \[
    A_B^2=(0,0,-1)^T,\qquad
    d_{\mathrm{ns},2}B^2_{02,20}
      =
    -\mathfrak b^2_{02,20},
    \qquad
    \widehat\sigma_{\mathrm{sc},2}(B^2_{02,20})=0,
  \]
  with the required locality or wavefront admissibility.  If such
  \(B^2_{02,20}\) is supplied, then
  \[
    d_{\mathrm{ns},2}(D^2_{02,20}+B^2_{02,20})
      =
    -2E^2_{\nu_{02}}+2E^2_{\nu_{20}}
      =
    -r_2 .
  \]

  For a compatible finite-window tower, write
  \(\mathfrak b^M\) for the corresponding Bianchi defect in window \(M\).
  The transport obstruction is
  \[
    \Delta^1_{M,N}
      :=
    -\pi_{M,N}\mathfrak b^M+\mathfrak b^N .
  \]
  A family \(B^M\) transports the fixed-window repair only when
  \[
    \Delta^1_{M,N}
      =
    d_{\mathrm{ns},N}(\pi_{M,N}B^M-B^N),
  \]
  possibly after an explicitly supplied scalar-zero transport correction.
  Thus the \(H^1\) class of \(\Delta^1_{M,N}\) must vanish first, and the
  resulting primitive-compatibility class in
  \[
    \varprojlim\nolimits^1_N
    H^0(\mathcal K^\bullet_{\mathrm{ns},N}(I))
  \]
  must vanish next.  Without these two gates, the lower \(\theta_3\)
  Bianchi repair is only a fixed-window matrix target, not a solved
  Costello/QME counterterm tower.
\end{prop}

\begin{proof}
  The source differential gives the two lower source-face rows with
  coefficients \(2\) and \(-2\).  Applying \(K^2_{\Theta_3}\) is not a
  chain map until its Hom-valued Bianchi defect is killed, so
  \[
    d_{\mathrm{ns},2}D^2_{02,20}
      =
    -2E^2_{\nu_{02}}+2E^2_{\nu_{20}}
      +\mathfrak b^2_{02,20},
  \]
  which is the column \(A_D^2=(-2,2,1)^T\) in the displayed basis.  The
  residual transported from the lower two-face row is \(r_2=(2,-2,0)^T\).
  If \(A_D^2c=-r_2=(-2,2,0)^T\), the first two coordinates force
  \(c=1\), while the Bianchi coordinate forces \(c=0\).  Thus the equation
  is inconsistent.  The same conclusion is the cokernel calculation:
  \(\widetilde\lambda_{02,\mathfrak b}\) annihilates \(A_D^2\) and does
  not annihilate \(r_2\).

  The target \((0,0,-1)^T\) is the boundary column required to remove
  only the Bianchi coordinate.  Since the same cokernel evaluates to
  \(-1\) on it, the current habitat does not contain such a boundary.
  Adding a scalar-zero local functional \(B^2_{02,20}\) with the displayed
  boundary produces
  \(A_D^2+A_B^2=(-2,2,0)^T=-r_2\), hence the lower residual is killed.
  This is only a formal finite-window repair until the local functional
  and its admissibility are supplied.

  For \(M\ge N\), applying \(d_{\mathrm{ns},N}\) to
  \(\pi_{M,N}B^M-B^N\) gives the difference between the transported
  Bianchi boundary and the \(N\)-window Bianchi boundary, namely
  \(-\pi_{M,N}\mathfrak b^M+\mathfrak b^N\).  Therefore the transported
  family exists only after this degree-one class vanishes in
  \(H^1(\mathcal K^\bullet_{\mathrm{ns},N}(I))\); after choosing
  primitives, their failure to be compatible is the usual Roos
  \(\varprojlim^1H^0\) obstruction.
\end{proof}

\begin{defn}[Normal \(\Omega\)-equivariant Costello kernel datum]
\label{def:app-normal-omega-costello-kernel-datum}
  Let
  \[
    X=\R_t\times\R_s\times\C_{z_1}\times\C_{z_2},
    \qquad
    L=\R_t\times\{s=z_1=z_2=0\},
  \]
  and let the normal torus be
  \[
    T_\Omega=
    \C^*_{\epsilon_s}\times
    \C^*_{\epsilon_1}\times
    \C^*_{\epsilon_2}.
  \]
  Its infinitesimal vector field is
  \[
    V_\Omega=
      \epsilon_s s\partial_s
      +\epsilon_1z_1\partial_{z_1}
      +\epsilon_2z_2\partial_{z_2},
    \qquad V_\Omega(t)=0,
  \]
  and the Cartan differential is
  \[
    Q_\Omega=Q+\iota_{V_\Omega},
    \qquad
    Q_\Omega^2=L_{V_\Omega}.
  \]
  All complexes below are taken in the \(T_\Omega\)-invariant, or
  equivalently basic Cartan, subcomplex.  Before this restriction the
  displayed operator is not a differential.

  Put
  \[
    \mathsf X_{N,M}
      \subset
    \{\,n\epsilon_s+a\epsilon_1+b\epsilon_2,\,
        n\epsilon_s-a\epsilon_1-b\epsilon_2\,\}
  \]
  for the finite set of nonzero moving normal characters occurring in the
  chosen \(N\)-brane, \(M\)-window package.  The coefficient ring is
  \[
    R_\Omega^{N,M}
      =
    \C[\epsilon_s,\epsilon_1,\epsilon_2,\hbar_\omega]
      [\chi^{-1}\mid\chi\in\mathsf X_{N,M}],
    \qquad
    \operatorname{wt}(\hbar_\omega)=
      \epsilon_1+\epsilon_2 .
  \]
  The self-dual branch is formed by base change to
  \(\epsilon_1+\epsilon_2=0\) before localization; only the surviving
  nonzero characters in the image of \(\mathsf X_{N,M}\) are inverted.
  In particular \(\epsilon_1+\epsilon_2\) is never inverted and then
  specialized to zero.  On that branch \(\hbar_\omega\) has weight zero
  and may be retained as a scalar parameter or set to \(1\) after the
  specialization is explicit.

  The three quantum parameters are distinct:
  \[
    \hbar_{\mathrm{QME}},\qquad
    \hbar_{\mathrm{W}},\qquad
    \hbar_\omega .
  \]
  The first counts Costello/BV graph order; the second deforms the
  boundary Weyl/Moyal algebra; the third scalarizes the equivariant
  Hamiltonian bracket.  A Weyl branch may identify
  \(\hbar_{\mathrm{W}}=\hbar_\omega\).  No branch identifies
  \(\hbar_{\mathrm{QME}}\) with either parameter unless a theorem states
  that extra specialization.  The equivariant Hamiltonian bracket is the
  \(R_\Omega^{N,M}\)-linear bracket
  \[
    \{f,g\}_\Omega
      =
    \hbar_\omega
    \bigl(
      \partial_{z_1}f\,\partial_{z_2}g
      -
      \partial_{z_2}f\,\partial_{z_1}g
    \bigr).
    \label{eq:app-omega-weighted-hamiltonian-bracket}
  \]
  A normal \(\Omega\)-equivariant finite-window Costello kernel datum over
  \((I,w,M,\mathcal G_M)\) is a finite-window marked Costello package as in
  Definitions~\ref{defn:finite-window-graph-qme-package}
  and~\ref{def:app-finite-window-marked-costello-list}, together with the
  following additional data.
  \begin{enumerate}[label={(\(\Omega\arabic*\))}]
    \item A \(T_\Omega\)-locally finite weighted field complex
      \(\mathcal E^{M}_{w,\Omega}|_I\), source CE complex, local
      interaction \(I^w_{0,M,\Omega}\), and invariant local-functional
      target
      \[
        \mathfrak Q^\bullet_{\mathcal G,\Omega,M}(I)
        \subset
        \mathcal O^{\mathrm{bd}}_{\mathrm{loc}}
          (\mathcal E^{M}_{w,\Omega}|_I;
           A_{\partial,\Omega,\hbar_{\mathrm{W}}}^{\mathrm{pp},w,M}(I)
           \widehat\otimes R_\Omega^{N,M}
           [[\hbar_{\mathrm{QME}}]])
      \]
      with differential
      \[
        d_{M,\Omega}
          =
        Q_\Omega+
        \{I^w_{0,M,\Omega},-\}_{\hbar_{\mathrm{QME}},M,\Omega}
      \]
      on the invariant subcomplex.  The source CE differential
      \(d_{\mathrm{CE},\Omega}\) is formed from \(Q_\Omega\) and
      the bracket~\eqref{eq:app-omega-weighted-hamiltonian-bracket}.

    \item For every localized normal character
      \(\chi\in\mathsf X_{N,M}\), or for every surviving nonzero image of
      such a character on the self-dual branch, a normal contraction
      \(h_\chi\) satisfying
      \[
        Q_\Omega h_\chi+h_\chi Q_\Omega
          =
        \operatorname{id}-\operatorname{pr}_{\chi=0}
      \]
      on the corresponding \(\chi\)-weight summand, after inverting
      \(\chi\).  This is the normal-mode acyclicity datum.  It is not a
      QME counterterm and does not decide any graph curvature class.

    \item An equivariant finite-scale gauge fixing and propagator
      \[
        P^{w,M}_{\epsilon,L,\Omega}
          =
        \int_\epsilon^L
          (Q^{\mathrm{GF}}_{\Omega}\otimes 1)K^{w,M}_{t,\Omega}\,dt
      \]
      such that
      \[
        [Q_\Omega,P^{w,M}_{\epsilon,L,\Omega}]
          =
        K^{w,M}_{\epsilon,\Omega}-K^{w,M}_{L,\Omega},
        \qquad
        [L_{V_\Omega},P^{w,M}_{\epsilon,L,\Omega}]=0,
      \]
      and such that the coefficient Casimir, wavefront conditions, and
      finite-window weighted bounds required by the nonequivariant package
      remain valid over \(R_\Omega^{N,M}[[\hbar_{\mathrm{QME}}]]\).

    \item A named normal-localization normalization
      \(\nu_\Omega\).  The residue branch has
      \[
        \nu_\Omega^{\mathrm{res}}=1.
      \]
      The Euler branch has
      \[
        \nu_\Omega^{\mathrm{Eu}}
          =
        \sigma_s(\epsilon_s\epsilon_1\epsilon_2)^{-1},
      \]
      after the corresponding Euler denominators have been admitted in the
      chosen localized ring.  Scalar-contact rows, graph weights, and trace
      normalizations are always read with the selected branch.  The residue
      branch, Euler branch, and Euler-rescaled residue branch are not
      silently identified; a comparison map must state the real orientation
      sign \(\sigma_s\) and every denominator.

    \item Equivariant local counterterm representatives
      \[
        C^{M}_{\Gamma,w,\Omega}(\epsilon;B^\Gamma_\bullet)
        \in
        \mathfrak Q^0_{\mathcal G,\Omega,M}(I)
      \]
      which are \(T_\Omega\)-invariant, local on the relevant bulk
      diagonal or brane-collision stratum, compatible with
      finite-window truncation and admissible weight transport, and
      whose scalar-contact associated graded is computed using the same
      \(\nu_\Omega\).

    \item An associated-graded scalar-symbol extractor candidate
      \[
        \sigma_{\mathrm{sc},\Omega,M}\colon
        \operatorname{gr}_{F}
        \mathfrak Q^\bullet_{\mathcal G,\Omega,M}(I)
        \longrightarrow
        S^\bullet_{\Omega,M},
        \qquad
        S^\bullet_{\Omega,M}
          =
        C^\bullet_{\mathrm{Lie}}
          (\bar A;\C\gamma_{\mathbf 1})[1]
          [[\hbar_{\mathrm{QME}}]]
          \widehat\otimes R_\Omega^{N,M},
      \]
      which sends the two-Hamiltonian scalar contact to
      \[
        \nu_\Omega\,
        \hbar_{\mathrm{QME}}\,
        \operatorname{Str}_{\mathrm{cyc}}(\mathfrak m)\,
        [\{f,g\}_\Omega]_0\,\gamma_{\mathbf 1}.
        \label{eq:app-omega-scalar-contact-symbol}
      \]
      A filtered chain lift
      \(\widehat\sigma_{\mathrm{sc},\Omega,M}\) is not part of the datum
      until the obstruction test below has been passed.

    \item A degree-zero Hom-valued graph kernel
      \[
        \kappa^M_{\mathcal G,\Omega,w,I}
        \in
        \operatorname{Hom}_{\mathrm{cont}}
        \left(
          C^\bullet_{\mathrm{CE},\Omega,\mathrm{lab}}(\mathcal G_M),
          \mathfrak Q^\bullet_{\mathcal G,\Omega,M}(I)
        \right)
      \]
      whose zero-edge part is the coefficient-current kernel and whose
      positive graph components are the renormalized
      \(\Omega\)-equivariant graph weights.
  \end{enumerate}
  This datum gives a habitat in which the equivariant QME and centrality
  obstructions are defined.  Localization supplies the coefficient ring,
  \(Q_\Omega\), normal contractions, and equivariant propagator data; it is
  not a QME proof and does not assert that any obstruction vanishes.
\end{defn}

\begin{thm}[Equivariant brane-defect QME obstruction criterion]
\label{thm:app-equivariant-brane-defect-qme-criterion}
  Fix a candidate normal \(\Omega\)-equivariant finite-window package
  satisfying items \((\Omega1)\)--\((\Omega6)\) of
  Definition~\ref{def:app-normal-omega-costello-kernel-datum}.  Let
  \[
    \sigma_{\mathrm{sc},\Omega,M}\colon
    \operatorname{gr}_{F}\mathfrak Q^\bullet_{\mathcal G,\Omega,M}(I)
    \longrightarrow
    S^\bullet_{\Omega,M}
  \]
  be its associated-graded scalar extractor.  The scalar projection is a
  theorem-grade chain projection exactly when the following relative Hom
  obstructions vanish.  First,
  \[
    \delta^{(0)}_{\Omega,\sigma,M}
      =
    d_{\mathrm{CE},\Omega}\sigma_{\mathrm{sc},\Omega,M}
      -
    \sigma_{\mathrm{sc},\Omega,M}d_{\operatorname{gr},\Omega}
      =
    0 .
  \]
  If this holds and \(s_{\Omega,M}\) is any filtered lift of
  \(\sigma_{\mathrm{sc},\Omega,M}\), the first nonzero homogeneous piece
  of
  \[
    d_{\mathrm{CE},\Omega}s_{\Omega,M}
      -
    s_{\Omega,M}d_{M,\Omega}
  \]
  lowering the scalar-contact filtration by \(r>0\) gives
  \[
    \mathfrak o^{(r)}_{\Omega,\sigma,M}
      \in
    H^1\!\left(
      \operatorname{Hom}_{T_\Omega,-r}
      \bigl(
        \operatorname{gr}_{F}
          \mathfrak Q^\bullet_{\mathcal G,\Omega,M}(I),
        S^\bullet_{\Omega,M}
      \bigr)
    \right).
  \]
  A filtered chain projection
  \(\widehat\sigma_{\mathrm{sc},\Omega,M}\) exists if and only if
  \(\delta^{(0)}_{\Omega,\sigma,M}=0\) and the classes
  \(\mathfrak o^{(r)}_{\Omega,\sigma,M}\) vanish successively, with
  compatible finite-window and weight-transport choices.

  Once this projection is supplied, put
  \[
    \mathcal K^\bullet_{\mathrm{ns},\Omega,M}(I)
      =
    \ker\widehat\sigma_{\mathrm{sc},\Omega,M},
    \qquad
    d_{\mathrm{ns},\Omega,M}
      =
    d_{M,\Omega}|_{\mathcal K_{\mathrm{ns},\Omega,M}} .
  \]
  Suppose the Hom-valued kernel of item \((\Omega7)\) is supplied, and
  counterterms have been chosen through order
  \(\hbar_{\mathrm{QME}}^{\,n-1}\).  The
  order-\(n\) equivariant residual is
  \[
  \begin{aligned}
    R^{\mathrm{ns}}_{n,\Omega,M}
      =
    [\hbar_{\mathrm{QME}}^n]\biggl(
      &\operatorname{Curv}_{\mathbf K,\Omega,M}
        (\kappa^M_{\mathcal G,\Omega,w,I})
       +d_{M,\Omega}C_{<n,\Omega,M} \\
      &\quad
       +\frac12
        \{C_{<n,\Omega,M},C_{<n,\Omega,M}\}
          _{\hbar_{\mathrm{QME}},M,\Omega}
    \biggr).
  \end{aligned}
  \]
  The scalar gate at this order is
  \[
    \mathfrak s_{n,\Omega,M}
      =
    \widehat\sigma_{\mathrm{sc},\Omega,M}
      (R^{\mathrm{ns}}_{n,\Omega,M}).
  \]
  Only when \(\mathfrak s_{n,\Omega,M}=0\) and all lower orders have been
  solved does the residual define the non-scalar obstruction class
  \[
    \mathfrak o^{\mathrm{ns}}_{n,\Omega,M}
      =
    [R^{\mathrm{ns}}_{n,\Omega,M}]
      \in
    H^1\!\left(
      \mathcal K^\bullet_{\mathrm{ns},\Omega,M}(I),
      d_{\mathrm{ns},\Omega,M}
    \right).
  \]
  A compatible order-\(n\) counterterm exists precisely when
  \[
    \bigl(\mathfrak o^{\mathrm{ns}}_{n,\Omega,M}\bigr)_M=0
    \quad\text{in}\quad
    \varprojlim_M
    H^1(\mathcal K^\bullet_{\mathrm{ns},\Omega,M}(I))
  \]
  and the Milnor primitive-compatibility class
  \[
    \lambda_{n,\Omega}
      \in
    \varprojlim\nolimits^1_M
    H^0(\mathcal K^\bullet_{\mathrm{ns},\Omega,M}(I))
  \]
  vanishes.  Iterating this criterion for all \(n\) is equivalent to
  \[
    \operatorname{Curv}_{\mathbf K,\Omega}
      (\kappa_{\mathcal G,\Omega,w,I}+C_{\Omega,w})
      =0
  \]
  in the completed invariant kernel complex.

  Therefore the normal \(\Omega\)-background supplies only the habitat.
  A QME proof requires, in addition, the chain scalar projection, the
  non-scalar row classes displayed above, explicit local counterterms
  killing those classes, compatible primitive transport in the
  \(\varprojlim^1H^0\) sense, and the \(Q_\Omega\)-centrality homotopies of
  Proposition~\ref{prop:app-qomega-centrality-homotopy-criterion} for the
  source rows whose centrality is used.
\end{thm}

\begin{proof}
  The first paragraph is the equivariant version of
  Lemma~\ref{lem:filtered-scalar-projection-obstruction}.  On the
  \(T_\Omega\)-invariant subcomplex \(Q_\Omega^2=L_{V_\Omega}=0\), so the
  Hom differential is
  \(D_\Omega(\alpha)=d_{\mathrm{CE},\Omega}\alpha
  -(-1)^{|\alpha|}\alpha d_{\operatorname{gr},\Omega}\).  The
  associated-graded defect is the obstruction to even starting a filtered
  lift.  After it vanishes, the first filtration-lowering chain defect is
  closed for \(D_\Omega\); changing the lift by a filtration-lowering
  cochain changes it by a Hom coboundary.  Completeness gives the
  successive criterion and the compatibility condition in the inverse
  finite-window tower.

  With a chain scalar projection fixed, its kernel is preserved by
  \(d_{M,\Omega}\).  The displayed residual is the order-\(n\) term of the
  curved Maurer--Cartan equation for the Hom-valued kernel plus lower
  counterterms.  Applying the scalar projection gives the scalar gate.  If
  the gate vanishes and the lower-order equations hold, the equivariant
  Bianchi identity in the invariant complex makes
  \(R^{\mathrm{ns}}_{n,\Omega,M}\) a
  \(d_{\mathrm{ns},\Omega,M}\)-cocycle.  Adding an order-\(n\) local
  counterterm changes this cocycle by a
  \(d_{\mathrm{ns},\Omega,M}\)-boundary, so the fixed-window obstruction
  is its \(H^1\)-class.  Compatibility of primitives across the
  finite-window inverse system is exactly the usual
  \(\varprojlim^1H^0\) obstruction of
  Theorem~\ref{thm:wt-nonscalar-counterterm-criterion}.  The final
  sentence follows because the normal contraction datum appears only in
  the construction of the equivariant propagator and localized coefficient
  complex; it does not identify the curvature residual with a boundary,
  does not produce scalar-projection chain data, and does not supply the
  centrality homotopies required for source rows.
\end{proof}

\begin{thm}[Minimal full-equivariant kernel vanishing]
\label{thm:app-minimal-full-equivariant-kernel-vanishing}
  Work in the local
  \(\R^2_{\mathrm{top}}\times\widehat{\C^2}_{0,\mathrm{hol}}\)
  full-equivariant marked branch with \(V=\C^{N|N}\).  Let the finite
  marked list be the minimal codimension-two closed list generated by the
  strict current-interface rows and by the one-cross two-Hamiltonian
  scalar-contact row of
  Proposition~\ref{prop:first-finite-window-array-rows}.  Assume that no
  further \(|\Gamma|_\hbar=1\) graph row, no genuine
  \(|\Gamma|_\hbar=2\) seed graph, no order-\(2\) CE-source face, and no
  order-\(2\) graph counterterm face is added.

  Then the first kernel residuals are
  \[
    R^{\mathrm{ns}}_{1,M}=0,\qquad
    R^{\mathrm{ns}}_{2,M}=0,
  \]
  and the compatible counterterms are
  \[
    C_{1,M}=0,\qquad C_{2,M}=0 .
  \]
  Thus the minimal full-equivariant finite-window array has no genuinely
  non-scalar QME obstruction through the first codimension-two closure
  level.  Its first possible non-scalar obstruction can only appear after
  a larger finite marked list supplies an additional scalar-zero row.
\end{thm}

\begin{proof}
  The order-\(0\) current-interface rows have zero curvature by
  Theorem~\ref{thm:app-algebraic-factorization-centre-zero-homotopies} and
  by the finite-window compatibility recorded in
  Proposition~\ref{prop:app-finite-window-marked-habitat-compatibility}.
  The only order-\(1\) row in the minimal list is the one-cross
  scalar-contact row.  Its scalar coefficient is the cyclic supertrace of
  the transported full-equivariant marker tensor.  By
  Lemma~\ref{lem:app-full-equivariant-cyclic-supertrace}, this coefficient
  is zero for \(V=\C^{N|N}\), including the empty word.  Hence the row is
  zero as a local functional in the full-equivariant branch, so
  \(R^{\mathrm{ns}}_{1,M}=0\) and \(C_{1,M}=0\).

  Codimension-two closure adds only the two ordered composites of two
  first scalar-contact boundary faces.  Their incidence signs are opposite
  and their output graphs, coefficients, and marker transports are
  identified by Equation~\eqref{eq:app-codim-two-marked-closure}.  This
  gives the pairwise cancellation of the marked differential.  In the
  full-equivariant branch each scalar-loop factor is already zero by the
  same cyclic supertrace calculation.  Since no genuine order-\(2\) seed graph,
  CE-source face, or counterterm face is present, there is no remaining
  row in the kernel complex.  Thus \(R^{\mathrm{ns}}_{2,M}=0\), and the
  zero counterterm is compatible under every finite-window truncation.
\end{proof}

\begin{prop}[First missing non-scalar row]
\label{prop:app-first-missing-nonscalar-row}
  Let \(\mathcal G^{\mathrm{mk}}_M\) be any larger finite marked list in
  the same full-equivariant local branch, with scalar projection killed as
  in Theorem~\ref{thm:app-full-equivariant-marked-shadow-vanishing}.  The
  first genuinely non-scalar obstruction not detected by the scalar
  projection is the least order \(n\) for which the row sum of
  Definition~\ref{def:app-nonscalar-kernel-row-complex} is nonzero in
  cohomology:
  \[
    \mathfrak o^{\mathrm{ns}}_{n,M}
      =
    \left[
      \sum_{\alpha\in\mathsf A^{\mathrm{fw}}_{n,M}}
      R^{\mathrm{row}}_{\alpha,M}
    \right]
    \in H^1\!\left(
      \mathcal K^\bullet_{\mathrm{ns},M}(I),d_{\mathrm{ns},M}
    \right).
  \]
  Since the minimal order-\(1\) scalar-contact row vanishes, the first
  datum that can make \(n=1\) is an additional order-\(1\) scalar-zero
  row
  \[
    \alpha\in
    \mathsf A^{\mathrm{fw}}_{1,M}\setminus\{\alpha_{\mathrm{sc}}\},
    \qquad
    S_{\alpha,M}=0,
  \]
  with its actual row value \(R^{\mathrm{row}}_{\alpha,M}\).  Equivalently,
  the missing row entry is one of the order-\(1\) summands
  \[
    \varepsilon_\Gamma d_M K^M_\Gamma,\qquad
    -\varepsilon^{\mathrm{CE}}_{\Gamma'}K^M_{\Gamma'},\qquad
    \frac12\varepsilon_{\Gamma_0,\Gamma_1}
      \{K^M_{\Gamma_0},K^M_{\Gamma_1}\}_{\hbar,M},
  \]
  together with its graph type, finite-window labels, full-equivariant
  marker tensor, incidence sign, scalar gate, and truncation coefficients
  \(t^{M,N}_{\alpha\alpha'}\).

  If the finite marked list explicitly has no such order-\(1\) row, the
  next missing datum is the order-\(2\) scalar-zero row sum
  \[
  \begin{aligned}
    R^{\mathrm{ns}}_{2,M}
      ={}&
      \sum_{|\Gamma|_\hbar=2}
        \varepsilon_\Gamma\,d_M K^M_\Gamma
      -\sum_{|\Gamma'|_\hbar=2}
        \varepsilon^{\mathrm{CE}}_{\Gamma'}K^M_{\Gamma'}\\
      &\quad
      +\frac12
      \sum_{|\Gamma_1|_\hbar=|\Gamma_2|_\hbar=1}
        \varepsilon_{\Gamma_1,\Gamma_2}
        \{K^M_{\Gamma_1},K^M_{\Gamma_2}\}_{\hbar,M}
      +[\hbar^2]d_M C_{<2,M}\\
      &\quad
      +\frac12[\hbar^2]
        \{C_{<2,M},C_{<2,M}\}_{\hbar,M}.
  \end{aligned}
  \]
  Without these row values, signs, scalar gates, and truncation
  coefficients, no theorem-grade non-scalar obstruction class or
  counterterm is defined beyond the minimal vanishing theorem above.
\end{prop}

\begin{proof}
  The scalar projection is a chain map on the full-equivariant marked
  habitat by Theorem~\ref{thm:app-full-equivariant-marked-shadow-vanishing},
  so its kernel is a \(d_M\)-stable complex.  After the scalar gate
  vanishes, Theorem~\ref{thm:finite-window-graph-qme-assembly} puts the
  first unsolved residual in degree one of that kernel complex.  Adding a
  degree-zero counterterm changes the residual by a \(d_M\)-boundary, and
  hence the displayed cohomology class is exactly the obstruction.

  The minimal order-\(1\) row is already zero by
  Theorem~\ref{thm:app-minimal-full-equivariant-kernel-vanishing}.  Thus an
  order-\(1\) non-scalar class can only be supplied by an additional
  finite-window row outside \(\alpha_{\mathrm{sc}}\).  The three displayed
  forms are precisely the order-\(1\) possibilities in the curvature
  expansion: target differential of an order-\(1\) graph component, CE
  source face of such a component, or a convolution bracket whose orders
  add to \(1\).  Each is a row only after the finite graph list supplies
  its labels, signs, marker transport, scalar gate, and truncation
  coefficients.

  If there is no order-\(1\) row, the next possible residual is the
  order-\(2\) curvature expansion, including genuine order-\(2\) graph
  components, CE-source faces, brackets of order-\(1\) components, and
  lower-counterterm terms.  These data are not determined by the scalar
  shadow theorem.  They are exactly the row entries required by
  Definition~\ref{defn:finite-window-graph-array}.
\end{proof}

\begin{prop}[First scalar-zero order-one row outside the minimal array]
\label{prop:app-first-scalar-zero-order-one-row}
  Work in the full-equivariant local branch of
  Theorem~\ref{thm:app-minimal-full-equivariant-kernel-vanishing}.  Let
  \(M\geq1\), and enlarge the minimal finite marked list by the labelled
  one-cross scalar-contact boundary operation
  \(\alpha_{11}\) with ordered Hamiltonian inputs
  \((a,z_1)\), \((b,z_1)\), the boundary Weyl edge \(E_{\mathrm{Weyl}}\),
  the scalar-contact output vertex \(\gamma_{\mathbf 1}\), and the same
  full-equivariant marker habitat as before.  Choose the orientation in
  which the displayed ordered pair has incidence coefficient \(+1\).

  Then \(\alpha_{11}\) is a scalar-zero order-one row outside
  \(\alpha_{\mathrm{sc}}\), with row data
  \[
    |\alpha_{11}|_\hbar=1,\qquad
    d_{\max}(\alpha_{11})=1,\qquad
    p_{\alpha_{11}}=1,\qquad
    r_{\alpha_{11}}=1,\qquad
    \iota_{\alpha_{11},M}=+1.
  \]
  Its row value and scalar gate are
  \[
    R^{\mathrm{row}}_{\alpha_{11},M}
      =
    \operatorname{Str}_{\mathrm{cyc}}
      \bigl(\mu_{\beta_{11}}(\mathfrak m)\bigr)
    \omega(z_1,z_1)
    \int_I a(t)b(t)\gamma_{\mathbf 1}(t)\,dt
      =0,
  \]
  \[
    S_{\alpha_{11},M}
      =
    \operatorname{Str}_{\mathrm{cyc}}
      \bigl(\mu_{\beta_{11}}(\mathfrak m)\bigr)
    \omega(z_1,z_1)\gamma_{\mathbf 1}=0.
  \]
  Here the coefficient is the extracted \(\hbar^1\)-coefficient; the full
  scalar-contact representative has the additional factor \(\hbar\).
  For \(M\geq N\geq1\), retain the same labelled row in window \(N\) and
  set
  \[
    t^{M,N}_{\alpha_{11}\alpha_{11}}=1,\qquad
    t^{M,N}_{\alpha_{11}\alpha'}=0
    \quad(\alpha'\neq\alpha_{11}).
  \]
  These coefficients satisfy the truncation equations because
  \(\pi_{M,N}R^{\mathrm{row}}_{\alpha_{11},M}=0\) and
  \(\pi_{M,N}S_{\alpha_{11},M}=0\).

  The associated kernel class is therefore
  \[
    \mathfrak o^{\mathrm{ns}}_{1,M}(\alpha_{11})
      =
    [R^{\mathrm{row}}_{\alpha_{11},M}]
      =
    0
    \in H^1\!\left(
      \mathcal K^\bullet_{\mathrm{ns},M}(I),d_{\mathrm{ns},M}
    \right),
  \]
  and the required counterterm contribution is \(C_{1,M}=0\).

  Thus the finite closed graph-list axioms do not support an absence
  theorem for literal scalar-zero order-one rows beyond
  \(\alpha_{\mathrm{sc}}\).  The first explicit such row is harmless:
  it is zero before cohomology.  If the row convention retains only
  nonzero summands, then \(\alpha_{11}\) is omitted, and a nonzero
  order-one class still requires an additional supplied row
  \(R^{\mathrm{row}}_{\alpha,M}\) of the form stated in
  Proposition~\ref{prop:app-first-missing-nonscalar-row}.
\end{prop}

\begin{proof}
  The scalar-contact boundary formula in
  Proposition~\ref{prop:first-finite-window-array-rows} applies to every
  ordered pair of Hamiltonian labels in the finite window.  For the pair
  \((z_1,z_1)\),
  \[
    \omega(z_1,z_1)
      =
    \bigl[\{z_1,z_1\}_{\Omega_{\mathrm{loc}}}\bigr]_0=0.
  \]
  Multiplying by the transported full-equivariant cyclic supertrace cannot
  change this zero coefficient.  Hence both the local functional row and
  its scalar projection vanish.

  Since the row value is zero, \(d_{\mathrm{ns},M}\)-closedness is
  automatic and the cohomology class is the zero class.  The displayed
  truncation coefficients send the labelled row to the same labelled row in
  lower windows containing \(z_1\), and both sides of the row and scalar
  truncation identities are zero.  The zero counterterm solves the row
  equation.  The final assertion is logical: the finite-list axioms allow
  this added labelled zero row, so they cannot prove that no scalar-zero
  order-one row exists; they can only decide a nonzero order-one class after
  the actual nonzero order-one row values have been supplied.
\end{proof}

\begin{rmk}[Order two after no nonzero order-one row]
\label{rmk:app-order-two-after-no-nonzero-order-one-row}
  If zero rows such as \(\alpha_{11}\) are suppressed and the finite marked
  list supplies no nonzero order-one row outside
  \(\alpha_{\mathrm{sc}}\), then \(C_{1,M}=0\) and every order-two bracket
  row containing \(\alpha_{\mathrm{sc}}\) is zero by
  Proposition~\ref{prop:universal-scalar-contact-bracket-rows}.  The
  order-two scalar-zero residual reduces to
  \[
  \begin{aligned}
    R^{\mathrm{ns}}_{2,M}
      ={}&
      \sum_{|\Gamma|_\hbar=2}
        \varepsilon_\Gamma\,d_M K^M_\Gamma
      -\sum_{|\Gamma'|_\hbar=2}
        \varepsilon^{\mathrm{CE}}_{\Gamma'}K^M_{\Gamma'}  \\
      &\quad
      +\frac12
      \sum_{\substack{|\Gamma_1|_\hbar=|\Gamma_2|_\hbar=1\\
        \Gamma_i\neq\alpha_{\mathrm{sc}}}}
        \varepsilon_{\Gamma_1,\Gamma_2}
        \bigl\{K^M_{\Gamma_1},K^M_{\Gamma_2}\bigr\}_{\hbar,M}.
  \end{aligned}
  \]
  Under the stronger minimal hypotheses of
  Theorem~\ref{thm:app-minimal-full-equivariant-kernel-vanishing} the
  three sums are empty, and \(R^{\mathrm{ns}}_{2,M}=0\).  In any larger
  finite package this displayed row sum is a cochain only after the
  genuine order-two graph weights, CE-source faces, bracket signs, scalar
  gates, and truncation coefficients have been supplied.
\end{rmk}

\begin{thm}[Minimal full-equivariant all-order vanishing]\label{thm:app-minimal-full-equivariant-all-order-vanishing}
  Work in the local
  \(\R^2_{\mathrm{top}}\times\widehat{\C^2}_{0,\mathrm{hol}}\)
  full-equivariant marked branch with \(V=\C^{N|N}\).  Let
  \(\mathcal G^{\mathrm{mk},0}_M\) be the finite marked list generated by
  the strict current-interface rows, the two-Hamiltonian scalar-contact
  boundary row \(\alpha_{\mathrm{sc}}\), and its subset-deletion closure in
  the sense of
  Definition~\ref{def:app-finite-window-marked-costello-list}.  Assume that no
  further positive-order seed graph, CE-source face, graph counterterm face,
  or nonzero scalar-zero row is added to the finite-window package.

  Then for every \(n\geq1\)
  \[
    R^{\mathrm{ns}}_{n,M}=0,\qquad
    \mathfrak o^{\mathrm{ns}}_{n,M}=0,\qquad
    C_{n,M}=0
  \]
  If the row convention suppresses zero rows, then
  \(\mathsf A^{\mathrm{fw},0}_{n,M}\) is empty for every \(n\geq1\).  If it
  retains labelled zero rows, every positive-order retained row has
  \[
    R^{\mathrm{row}}_{\alpha,M}=0,\qquad
    S_{\alpha,M}=0,
  \]
  and may be given the zero truncation row
  \(t^{M,N}_{\alpha\alpha'}=0\).  In particular, the minimal
  full-equivariant order-three residual is zero.
\end{thm}

\begin{proof}
  Proposition~\ref{prop:first-finite-window-array-rows} proves that the
  one-cross scalar-contact component satisfies
  \(K^M_{\alpha_{\mathrm{sc}}}=0\) as a local functional, not only after
  scalar projection.  The proof is the cyclic-supertrace calculation of
  Lemma~\ref{lem:app-full-equivariant-cyclic-supertrace}: every transported
  full-equivariant marker on a scalar open loop has cyclic supertrace zero,
  including the empty marker word.

  Every positive-order boundary row in
  \(\mathcal G^{\mathrm{mk},0}_M\) contains at least one such
  scalar-contact factor.  Hence its local-functional value is zero.  Bracket
  rows in the minimal package have at least one positive-order factor of this
  form, so bilinearity of the BV/convolution bracket gives zero.  Lower
  counterterm rows vanish inductively because the chosen lower counterterms
  are zero.  Genuine \(dK\), CE-source, and graph-counterterm rows are absent
  by the minimality hypothesis.  Therefore the order-\(n\) scalar-zero row sum
  is zero for all \(n\geq1\), its class in
  \(H^1(\mathcal K^\bullet_{\mathrm{ns},M}(I),d_{\mathrm{ns},M})\) is zero,
  and the zero counterterm is a compatible primitive.  Truncation sends zero
  rows to zero rows, giving the displayed truncation choice.
\end{proof}

\begin{prop}[First order-three larger-package datum]\label{prop:app-first-order-three-larger-package-datum}
  Keep the full-equivariant local branch, and suppose the minimal package has
  been closed through order two with
  \(C_{1,M}=C_{2,M}=0\) and no nonzero order-one or order-two row outside the
  scalar-contact closure.  Then the finite graph-list axioms do not by
  themselves supply a nonzero order-three obstruction.  A nonzero order-three
  cochain exists only after the finite-window package supplies at least one
  of the following row data:
  \[
  \begin{aligned}
    R^{\mathrm{ns}}_{3,M}
      ={}&
      \sum_{|\Gamma|_\hbar=3}
        \varepsilon_\Gamma\,d_M K^M_\Gamma
      -\sum_{|\Gamma'|_\hbar=3}
        \varepsilon^{\mathrm{CE}}_{\Gamma'}K^M_{\Gamma'}\\
      &\quad
      +\frac12
      \sum_{\substack{|\Gamma_1|_\hbar+|\Gamma_2|_\hbar=3\\
        K^M_{\Gamma_1}\neq0,\;K^M_{\Gamma_2}\neq0}}
        \varepsilon_{\Gamma_1,\Gamma_2}
        \{K^M_{\Gamma_1},K^M_{\Gamma_2}\}_{\hbar,M}\\
      &\quad
      +[\hbar^3]\left(
        d_M C_{<3,M}
        +\frac12\{C_{<3,M},C_{<3,M}\}_{\hbar,M}\right).
  \end{aligned}
  \]
  Under the preceding minimal hypotheses the bracket and lower-counterterm
  sums are empty.  Thus the first missing order-three datum is a genuine
  \(|\Gamma|_\hbar=3\) graph-differential row
  \[
    R^{\mathrm{row}}_{d\Gamma,M}
      =
    \varepsilon_\Gamma\,d_M K^M_\Gamma,
    \qquad
    S_{d\Gamma,M}
      =
    \varepsilon_\Gamma\,\widehat\sigma_{\mathrm{sc},M}(d_M K^M_\Gamma),
  \]
  or an order-three CE-source row
  \[
    R^{\mathrm{row}}_{\mathrm{CE}(\Gamma,\nu),M}
      =
    -\varepsilon^{\mathrm{CE}}_{\Gamma,\nu}
      K^M_\Gamma(\zeta_{M,\nu}),
    \qquad
    S_{\mathrm{CE}(\Gamma,\nu),M}
      =
    -\varepsilon^{\mathrm{CE}}_{\Gamma,\nu}
      \widehat\sigma_{\mathrm{sc},M}
      \bigl(K^M_\Gamma(\zeta_{M,\nu})\bigr).
  \]
  In a nonminimal package with supplied nonzero order-one and order-two
  components one must also supply each bracket row
  \[
    R^{\mathrm{row}}_{b(\Gamma_1,\Gamma_2),M}
      =
    \frac12\varepsilon_{\Gamma_1,\Gamma_2}
    \{K^M_{\Gamma_1},K^M_{\Gamma_2}\}_{\hbar,M},
    \qquad
    S_{b,M}
      =
    \widehat\sigma_{\mathrm{sc},M}
      (R^{\mathrm{row}}_{b(\Gamma_1,\Gamma_2),M}),
  \]
  for \(|\Gamma_1|_\hbar+|\Gamma_2|_\hbar=3\), together with its
  truncation coefficients \(t^{M,N}_{bb'}\).

  The scalar gate for order three is the single condition
  \[
    \mathfrak s_{3,M}
      =
    \sum_{\alpha\in\mathsf A^{\mathrm{fw}}_{3,M}}S_{\alpha,M}=0.
  \]
  Only after this gate is supplied and vanishes is the non-scalar obstruction
  class defined:
  \[
    \mathfrak o^{\mathrm{ns}}_{3,M}
      =
    [R^{\mathrm{ns}}_{3,M}]
    \in H^1\!\left(
      \mathcal K^\bullet_{\mathrm{ns},M}(I),d_{\mathrm{ns},M}
    \right).
  \]
  A fixed-window counterterm \(C_{3,M}\) exists exactly when this class
  vanishes; a compatible tower also requires the corresponding
  \(\varprojlim^1 H^0\)-compatibility class to vanish.
\end{prop}

\begin{proof}
  The order-three expansion is
  Theorem~\ref{thm:actual-finite-window-nonscalar-decision-datum} with
  \(n=3\), written in the row notation of
  Definition~\ref{defn:finite-window-graph-array}.  Rows containing
  \(\alpha_{\mathrm{sc}}\) vanish by
  Proposition~\ref{prop:universal-scalar-contact-bracket-rows}, and rows
  containing any retained labelled zero scalar-contact row vanish by the same
  bilinearity argument.  Since \(C_{1,M}=C_{2,M}=0\), the lower-counterterm
  contribution is zero in the minimal package.  Therefore no nonzero
  order-three cochain is produced unless the package supplies a genuine
  order-three graph row, an order-three CE-source row, or a bracket row between
  already supplied nonzero lower-order components.  The displayed formulas are
  precisely the row values and scalar gates required by
  Definition~\ref{defn:finite-window-graph-array}.  The final assertions are
  the kernel obstruction criterion of
  Definition~\ref{def:app-nonscalar-kernel-row-complex} and the finite-window
  compatibility criterion of
  Theorem~\ref{thm:actual-finite-window-nonscalar-decision-datum}.
\end{proof}

\begin{prop}[Marked-row test for unreduced centrality homotopies]
\label{prop:app-marked-row-centrality-homotopy-test}
  Work in a finite-window marked Costello list for the local mixed HT
  brane-defect theory, with labels and source faces as in
  Definition~\ref{def:app-finite-window-marked-costello-list}.  Let
  \(\kappa^M_{\mathcal G,w,I}\) be the Hom-valued finite-window kernel
  carried by this list.  For source generators
  \(x,y\in\{u_{a(t)dt\otimes f},c_{B,\rho}\}\), the marked row testing
  unreduced centrality is
  \[
    R^{\mathrm{cent}}_{x,y,M}
      =
    \operatorname{Curv}_{\mathbf K,M}
      (\kappa^M_{\mathcal G,w,I})(x\wedge y).
  \]
  In row notation this is the finite sum
  \[
  \begin{aligned}
    R^{\mathrm{cent}}_{x,y,M}
      ={}&
      \sum_{\Gamma\in\mathsf A_{x,y,M}^{d}}
        \varepsilon_\Gamma\,d_M K^M_\Gamma
      -\sum_{\Gamma\in\mathsf A_{x,y,M}^{\mathrm{CE}}}
        \varepsilon^{\mathrm{CE}}_\Gamma
        K^M_\Gamma(\zeta_{\Gamma,x,y})  \\
      &\quad
      +\frac12
      \sum_{(\Gamma_1,\Gamma_2)\in
        \mathsf A_{x,y,M}^{\mathrm{br}}}
        \varepsilon_{\Gamma_1,\Gamma_2}
        \{K^M_{\Gamma_1},K^M_{\Gamma_2}\}_{\hbar,M}
      +R^{\mathrm{ct}}_{x,y,M}.
  \end{aligned}
  \]
  Here the four summands are respectively the target differential rows,
  the CE-source rows, the convolution bracket rows, and the supplied
  counterterm rows of the marked table.  The scalar gate is
  \[
    S^{\mathrm{cent}}_{x,y,M}
      =
    \widehat\sigma_{\mathrm{sc},M}
      (R^{\mathrm{cent}}_{x,y,M}).
  \]
  Only after this gate is zero does the centrality row define a class
  \[
    \mathfrak o^{\mathrm{cent}}_{x,y,M}
      =
    [R^{\mathrm{cent}}_{x,y,M}]
      \in
    H^{|x|+|y|}
      (\mathcal K^\bullet_{\mathrm{ns},M}(I),d_{\mathrm{ns},M}).
  \]

  A marked finite-window package supplies an actual unreduced
  centrality homotopy for the pair \((x,y)\) exactly when it also
  supplies a degree \(|x|+|y|-1\) row
  \(H^M_{x,y}\in\mathcal K_{\mathrm{ns},M}^{|x|+|y|-1}(I)\) with
  \[
    R^{\mathrm{cent}}_{x,y,M}+d_M H^M_{x,y}=0,
  \]
  and these rows satisfy the truncation and weight identities of
  Theorem~\ref{thm:app-curved-kernel-centrality-homotopy-criterion}.  If
  such a row or compatibility identity is absent from the table, the
  missing datum is exactly the omitted centrality-homotopy row, not a
  consequence of scalar-shadow vanishing.

  In the minimal full-equivariant branch of
  Theorem~\ref{thm:app-minimal-full-equivariant-all-order-vanishing},
  every positive-order row contains the zero scalar-contact factor
  \(K^M_{\alpha_{\mathrm{sc}}}=0\), and no additional positive-order seed
  graph, CE-source face, graph counterterm face, or nonzero scalar-zero
  row is present.  Hence
  \[
    R^{\mathrm{cent}}_{x,y,M}=0,\qquad H^M_{x,y}=0
  \]
  is a valid homotopy inside that minimal finite marked habitat.  This
  does not construct the centrality homotopy for an enlarged unreduced
  Costello graph/counterterm target; any larger package must supply the
  displayed row values, scalar gate, primitive \(H^M_{x,y}\), and
  transition coefficients.
\end{prop}

\begin{proof}
  Proposition~\ref{prop:app-finite-window-marked-habitat-compatibility}
  identifies the marked differential with the finite-window QME
  differential and, for Hom-valued kernels, adds the source face
  \(-\kappa d_{\mathrm{CE}}\).  Therefore the binary curvature on
  \(x\wedge y\) is exactly the sum of the four row types displayed above:
  target differential, CE-source subtraction, convolution bracket, and
  counterterm contribution.  Applying the scalar projection gives the
  scalar gate.  If the gate vanishes, the row lies in the non-scalar
  kernel complex of
  Definition~\ref{def:app-nonscalar-kernel-row-complex}; its cohomology
  class is the obstruction to finding a fixed-window primitive.

  The equation
  \(R^{\mathrm{cent}}_{x,y,M}+d_M H^M_{x,y}=0\) is precisely the
  homotopy
  equation~\eqref{eq:app-centrality-homotopy-equation}.  The transition
  identities are necessary because the local kernel is the inverse limit
  over finite windows and admissible weights.  Thus a row not present in
  the marked table cannot be inferred from the scalar projection; it is
  the missing graph, source-face, counterterm, or homotopy datum.

  In the minimal full-equivariant package all positive-order scalar-contact
  rows are zero before cohomology by
  Theorem~\ref{thm:app-minimal-full-equivariant-all-order-vanishing}.
  Since no other positive-order row is allowed by the minimal hypothesis,
  the centrality row is zero and the zero primitive is compatible with
  all truncation maps.  The final assertion is the contrapositive of the
  row formula: once an enlarged package adds a nonzero graph, source,
  bracket, or counterterm row, the corresponding scalar gate and primitive
  must be supplied explicitly.
\end{proof}

\begin{prop}[\(Q_\Omega\)-centrality homotopy criterion]
\label{prop:app-qomega-centrality-homotopy-criterion}
  Let
  \((I,w,M,\mathcal G_M)\) carry a normal \(\Omega\)-equivariant
  Costello kernel datum in the sense of
  Definition~\ref{def:app-normal-omega-costello-kernel-datum}, and assume
  the scalar projection has passed the Hom-obstruction test of
  Theorem~\ref{thm:app-equivariant-brane-defect-qme-criterion}.  For
  \(T_\Omega\)-invariant source generators
  \[
    x,y\in\{u_{a(t)dt\otimes f},c_{B,\rho}\},
  \]
  define the equivariant centrality row
  \[
    R^{\mathrm{cent}}_{x,y,\Omega,M}
      =
    \operatorname{Curv}_{\mathbf K,\Omega,M}
      (\kappa^M_{\mathcal G,\Omega,w,I})(x\wedge y).
  \]
  Its finite-row expansion is
  \[
  \begin{aligned}
    R^{\mathrm{cent}}_{x,y,\Omega,M}
      ={}&
      \sum_{\Gamma\in\mathsf A_{x,y,\Omega,M}^{d}}
        \varepsilon_\Gamma\,d_{M,\Omega}K^M_{\Gamma,\Omega}
      -\sum_{\Gamma\in\mathsf A_{x,y,\Omega,M}^{\mathrm{CE}}}
        \varepsilon^{\mathrm{CE}}_\Gamma
        K^M_{\Gamma,\Omega}(\zeta_{\Gamma,x,y,\Omega})  \\
      &\quad
      +\frac12
      \sum_{(\Gamma_1,\Gamma_2)\in
        \mathsf A_{x,y,\Omega,M}^{\mathrm{br}}}
        \varepsilon_{\Gamma_1,\Gamma_2}
        \{K^M_{\Gamma_1,\Omega},
          K^M_{\Gamma_2,\Omega}\}_{\hbar_{\mathrm{QME}},M,\Omega}
      +R^{\mathrm{ct}}_{x,y,\Omega,M}.
  \end{aligned}
  \]
  The scalar gate is
  \[
    S^{\mathrm{cent}}_{x,y,\Omega,M}
      =
    \widehat\sigma_{\mathrm{sc},\Omega,M}
      (R^{\mathrm{cent}}_{x,y,\Omega,M}).
  \]
  Only after this gate vanishes does one obtain the obstruction class
  \[
    \mathfrak o^{\mathrm{cent}}_{x,y,\Omega,M}
      =
    [R^{\mathrm{cent}}_{x,y,\Omega,M}]
      \in
    H^{|x|+|y|}
      (\mathcal K^\bullet_{\mathrm{ns},\Omega,M}(I),
       d_{\mathrm{ns},\Omega,M}).
  \]

  A \(Q_\Omega\)-centrality homotopy for \((x,y)\) is exactly a
  \(T_\Omega\)-invariant row
  \[
    H^\Omega_{x,y,M}
      \in
    \mathcal K^{|x|+|y|-1}_{\mathrm{ns},\Omega,M}(I)
  \]
  satisfying
  \[
    R^{\mathrm{cent}}_{x,y,\Omega,M}
      +d_{M,\Omega}H^\Omega_{x,y,M}=0,
    \qquad
    L_{V_\Omega}H^\Omega_{x,y,M}=0,
    \label{eq:app-qomega-centrality-homotopy-equation}
  \]
  together with the same interval, finite-window, and admissible-weight
  transition identities as in
  Theorem~\ref{thm:app-curved-kernel-centrality-homotopy-criterion}.
  Equivalently, the inverse-system class
  \[
    \bigl(\mathfrak o^{\mathrm{cent}}_{x,y,\Omega,M}\bigr)_M
    \in
    \varprojlim_M
    H^{|x|+|y|}
      (\mathcal K^\bullet_{\mathrm{ns},\Omega,M}(I))
  \]
  and the corresponding
  \(\varprojlim^1 H^{|x|+|y|-1}\)-primitive compatibility class must both
  vanish.

  If one works before passing to \(T_\Omega\)-invariants, then
  \(Q_\Omega^2=L_{V_\Omega}\) gives a secondary obstruction
  \[
    \ell^{\mathrm{cent}}_{x,y,\Omega,M}
      =
    [L_{V_\Omega}H_{x,y,\Omega,M}]
      \in
    H^{|x|+|y|}
      (\mathcal K^\bullet_{\mathrm{ns},\Omega,M}(I),
       d_{\mathrm{ns},\Omega,M})
  \]
  for any chosen primitive.  This class must vanish, or the primitive
  must be replaced by a basic one, before the homotopy is a
  \(Q_\Omega\)-homotopy.
\end{prop}

\begin{proof}
  The row expansion is the Hom-valued curvature computation of
  Proposition~\ref{prop:app-marked-row-centrality-homotopy-test} with
  \(d_M\), the CE differential, the bracket, and the counterterm table
  replaced by their equivariant counterparts from
  Definition~\ref{def:app-normal-omega-costello-kernel-datum}.  The
  scalar gate is necessary because only its vanishing places the row in
  \(\mathcal K^\bullet_{\mathrm{ns},\Omega,M}\).  Once there, the
  equivariant Bianchi identity makes the row closed, so its cohomology
  class is the obstruction to a fixed-window primitive.

  The displayed equation is precisely the centrality homotopy equation in
  the invariant \(Q_\Omega\)-complex.  Since \(Q_\Omega^2=L_{V_\Omega}\),
  invariance of the primitive is part of the datum, not a consequence of
  scalar-shadow vanishing.  The inverse-limit assertion is the same
  primitive-compatibility argument as in
  Theorem~\ref{thm:app-curved-kernel-centrality-homotopy-criterion}.  If
  a primitive is chosen outside the basic subcomplex, applying
  \(d_{M,\Omega}\) a second time records \(L_{V_\Omega}H\); its
  cohomology class is exactly the secondary obstruction displayed above.
\end{proof}

\begin{prop}[Even-block marker obstruction]
\label{prop:app-even-block-marker-obstruction}
  Replace full \(\lie{gl}(V)\)-equivariance by the smaller even
  block-diagonal habitat
  \[
    \lie{gl}(V_{\bar0})\oplus\lie{gl}(V_{\bar1})
      \subset \lie{gl}(V),
    \qquad
    V_{\bar0}\cong V_{\bar1}\cong\C^N .
  \]
  Then a single fixed even marker is admissible exactly when
  \[
    T=\lambda_0P_{\bar0}+\lambda_1P_{\bar1},
  \]
  where \(P_{\bar0}\) and \(P_{\bar1}\) are the parity block
  projections.  Its scalar shadow is
  \[
    \operatorname{Str}(T)=N(\lambda_0-\lambda_1).
  \]
  For a two-Hamiltonian scalar contact the obstruction class is
  \[
    \hbar N(\lambda_0-\lambda_1)\,
    \omega(f,g)\gamma_{\mathbf 1},
    \qquad
    \hbar N(\lambda_0-\lambda_1)[\bar c]\,\gamma_{\mathbf 1}
    \text{ in cohomology}.
  \]
  Since \([\bar c]\neq0\) in the scalar-reduced source, this smaller
  habitat has zero \([\bar c]\)-shadow for single markers if and only if
  \(\lambda_0=\lambda_1\), i.e. the marker is scalar and the habitat
  has collapsed back to the full-equivariant single-marker case.  For
  higher marker tensors the strongest true positive statement is exactly
  the supertraceless-monodromy subhabitat of
  Theorem~\ref{thm:app-supertraceless-monodromy-projection}; outside it
  the obstruction is the displayed
  \(\operatorname{Str}_{\mathrm{cyc}}\)-weighted \([\bar c]\)-class.
\end{prop}

\begin{proof}
  The commutant of
  \(\lie{gl}(V_{\bar0})\oplus\lie{gl}(V_{\bar1})\) on
  \(V_{\bar0}\oplus V_{\bar1}\) is
  \(\C P_{\bar0}\oplus\C P_{\bar1}\), so the displayed form of \(T\) is
  necessary and sufficient.  The supertrace is
  \[
    \operatorname{Str}(\lambda_0P_{\bar0}+\lambda_1P_{\bar1})
      =
    \lambda_0\dim V_{\bar0}-\lambda_1\dim V_{\bar1}
      =
    N(\lambda_0-\lambda_1).
  \]
  Proposition~\ref{prop:app-first-nonidentity-loop-scalar-symbol} then
  gives the scalar CE representative and the cohomology class by
  multiplying \(\omega\), respectively \([\bar c]\), by this coefficient.
  The class \([\bar c]\) is nonzero by
  Theorem~\ref{thm:app-scalar-contact-qme-class}; hence the class
  vanishes for all two-Hamiltonian tests exactly when
  \(\lambda_0=\lambda_1\).  The higher tensor assertion is the same
  calculation with
  \(\operatorname{Str}(T)\) replaced by
  \(\operatorname{Str}_{\mathrm{cyc}}\) and with the preservation
  criterion supplied by
  Theorem~\ref{thm:app-supertraceless-monodromy-projection}.
\end{proof}

\begin{thm}[Ordinary scalar-reduced \(\lie{gl}_N\) obstruction criterion]
\label{thm:app-ordinary-glN-scalar-obstruction-criterion}
  Work in the scalar-reduced source of
  Definition~\ref{def:app-scalar-obstruction-complex} and keep the
  ordinary trace on the open \(\lie{gl}_N\) brane.  Let
  \(W_{\mathrm{ext}}\) be any additional formal-local closed-exchange or
  brane-defect graph contribution whose leading scalar associated
  graded is \(w_{\mathrm{ext}}\in C^2_{\mathrm{Lie}}(\bar A;\C)\).  A
  cancellation
  \[
    \operatorname{Ob}_{\mathrm{sc}}
      +W_{\mathrm{ext}}+d_{\mathrm{QME}}C=0
  \]
  with \(C\in\mathcal O_{\mathrm{loc}}^{\mathrm{bd}}\) can hold only if
  \[
    [w_{\mathrm{ext}}]=-\hbar N[\bar c]
      \qquad\text{in }H^2_{\mathrm{Lie}}(\bar A;\C)[[\hbar]].
  \]
  In particular, with \(W_{\mathrm{ext}}=0\), or with
  \(w_{\mathrm{ext}}\) exact, the ordinary scalar-reduced
  \(\lie{gl}_N\) brane-defect model has no scalar QME cancellation.
\end{thm}

\begin{proof}
  Taking associated graded with respect to the scalar-contact filtration
  sends the equation
  \[
    \operatorname{Ob}_{\mathrm{sc}}
      +W_{\mathrm{ext}}+d_{\mathrm{QME}}C=0
  \]
  to
  \[
    \hbar N\omega+w_{\mathrm{ext}}+\delta\bar\eta=0
      \qquad\text{in }C^2_{\mathrm{Lie}}(\bar A;\C),
  \]
  where \(\bar\eta\) is the linear functional represented by the
  leading scalar part of \(C\).  Passing to cohomology gives the stated
  necessary condition.  If \(W_{\mathrm{ext}}=0\), or if
  \(w_{\mathrm{ext}}\) is exact, this condition would force
  \(\hbar N[\bar c]=0\) in \(H^2_{\mathrm{Lie}}(\bar A;\C)[[\hbar]]\).
  For \(N>0\) and formal \(\hbar\), this
  contradicts Theorem~\ref{thm:app-scalar-contact-qme-class}, which
  identifies \([\bar c]\neq0\) on \(\bar A\).  Thus no internal local
  counterterm cancels the ordinary scalar-reduced obstruction.
\end{proof}
